% Candidate statements and complete proofs for the DEG_L paper.
% Checked by /root/semireg_t7, GPT-6 Astra Pro; inherited effort not exposed.
% This is an input fragment, not a standalone article.
% Source locators are written explicitly so no guessed bibliography keys are used.

\subsection{The actual surface and its deformation problem}

\begin{proposition}[S-01: Allowed source covers do not preserve an embedding premise]
\label{prop:semireg-cover-domain}
Let $f:Y\to A$ be a surface morphism obtained from the four rational
Hodge maps in 032, Proposition~3.1. If $\pi:Y'\to Y$ is a connected finite
\'etale cover of degree $d>1$, the pulled-back maps still satisfy its
nonzero-period criterion, but $f\pi$ is not a closed immersion.
Moreover $(f\pi)_*[Y']=d f_*[Y]$.
\end{proposition}
\begin{proof}
Pullback preserves rational Hodge maps and multiplies the integral of
a top-degree class by $d$. Every geometric point of $Y$ has $d$ distinct
preimages with the same image under $f\pi$, so this morphism is not
injective on geometric points. The last assertion is the proper
pushforward formula for a finite map. The deeper principal covers
used here are explicitly permitted in 032, Corollary~6.7 and \S7.
\end{proof}

\begin{proposition}[S-02: The exact conormal gate]
\label{prop:semireg-dual-gate}
For a smooth closed immersion $i:Y\hookrightarrow A$ of a surface in
an abelian $g$-fold, $g\geq3$, write $N=N_{Y/A}$ and
$W=H^0(A,\Omega_A^1)$. Bloch semiregularity is equivalent to
surjectivity of
\[
 H^1(A,\mathcal O_A)\otimes\bigwedge^3W
 \longrightarrow H^1(Y,N^\vee\otimes K_Y).
\]
The map is restriction followed by the natural morphism
$\bigwedge^3 i^*\Omega_A^1\to N^\vee\otimes K_Y$.
\end{proposition}
\begin{proof}
Bloch's map is
$H^1(Y,N)\to H^{g-1}(A,\Omega_A^{g-3})$.
Serre duality identifies its dual source with $H^1(A,\Omega_A^3)$
and its dual target with $H^1(Y,N^\vee\otimes K_Y)$.
The conormal exact sequence induces the displayed exterior-power
map, since the quotient $\Omega_Y^1$ has rank two.
This is the defining dual map in Buchweitz--Flenner \S8.1;
Proposition~8.2 identifies it with Bloch's map.
Finally $\Omega_A^3=\mathcal O_A\otimes\bigwedge^3W$,
and finite-dimensional duality gives the equivalence.
\end{proof}

\begin{proposition}[S-03: A nonzero period does not force semiregularity]
\label{prop:semireg-period-counterexample}
There is a smooth surface in a CM abelian threefold with a nonzero
product of two holomorphic and two antiholomorphic one-forms for
which the Bloch semiregularity kernel has dimension at least eight.
\end{proposition}
\begin{proof}
Let $E$ be a CM elliptic curve, $B=E^2$, and
$L=\mathcal O_E(3[0])\boxtimes\mathcal O_E(3[0])$.
A general $C\in|L|$ is smooth and connected. Since $L^2=18$ and
$K_B=\mathcal O_B$, adjunction gives $g(C)=10$ and
$N_{C/B}=K_C$. Put $A=B\times E$ and $Y=C\times E$.
Then $N_{Y/A}=K_C\boxtimes\mathcal O_E$, so K\"unneth gives
$h^1(Y,N)=1+10=11$. The divisor semiregularity target
$H^2(A,\mathcal O_A)$ has dimension three.
The curve $C$ dominates the first elliptic factor: otherwise its
ample divisor class would be a fibre class. Hence a nonzero one-form
from that factor restricts to a nonzero form $\alpha$ on $C$.
For a nonzero one-form $\beta$ on the last $E$, the integral
$\int_Y(\alpha\wedge\beta)\wedge\overline{\alpha\wedge\beta}$
is nonzero, by the product formula and positivity of the two
one-dimensional integrals.
\end{proof}

\begin{proposition}[S-04: Rigid generating sources]
\label{prop:semireg-rigid-source}
Let $Y$ be a smooth projective surface with $H^1(Y,T_Y)=0$ and
$i:Y\hookrightarrow A$ a smooth closed immersion. Let $G\subset A$
be the abelian subvariety generated by a translate of $i(Y)$.
Then the ambient obstruction map satisfies
\[
\ker(\operatorname{ob}_i)=
\operatorname{Lie}(A)\otimes
\ker\bigl(H^1(A,\mathcal O_A)\to H^1(Y,\mathcal O_Y)\bigr).
\]
In particular it is injective when $G=A$. If in addition $i$ is
Bloch semiregular, this is also the infinitesimal Hodge-locus kernel
of the entire class $i_*[Y]$.
\end{proposition}
\begin{proof}
Restriction on $H^1(T_A)$ is $\operatorname{id}_{\operatorname{Lie}(A)}
\otimes i^*$, since $T_A$ is trivial. The normal sequence and
$H^1(T_Y)=0$ make $H^1(i^*T_A)\to H^1(N)$ injective.
This proves the kernel formula. The map $\operatorname{Alb}(Y)\to A$
has image $G$, so $i^*$ on $H^1(\mathcal O)$ is injective precisely
when $G=A$. Buchweitz--Flenner, Corollary~4.12 and Proposition~8.2,
identify
\[
 \sigma_i(\operatorname{ob}_i(\kappa))
 =\kappa\mathbin{\lrcorner}i_*[Y].
\]
If $\sigma_i$ is injective, the kernels of the two sides agree.
\end{proof}

The actual 032 source is a compact smooth complex two-ball quotient.
Calabi--Vesentini, \emph{Sur les vari\'et\'es complexes compactes
localement sym\'etriques}, Th\'eor\`eme~3, p.~315, gives
$H^1(Y,T_Y)=0$. The nonzero four-factor period does not establish
$G=A$. For example, separate injectivity on the two eigenspaces
of a quadratic action is insufficient: in
$A=E^n\times E^n$ with action $(a,\bar a)$, pullback to the diagonal
$E^n$ is injective on each eigenspace, but their images coincide.

\begin{proposition}[S-05: A quantitative obstruction for the actual morphism]
\label{prop:semireg-map-dimension}
Let $Y$ be a compact smooth complex two-ball quotient and
$f:Y\to A$ a morphism with two-dimensional image, where $\dim A=g$.
Let $\tau_f$ be the Buchweitz--Flenner morphism semiregularity map.
Then
\[
 \dim\ker\tau_f\ \geq\ gq(Y)-\binom{2g}{g-2}.
\]
\end{proposition}
\begin{proof}
In characteristic zero $df:T_Y\to f^*T_A=\mathcal O_Y^{\oplus g}$
is generically injective. Its kernel is a torsion-free sheaf of
generic rank zero, hence vanishes. With $N_f=\operatorname{coker}(df)$,
the derived dual of $L_{Y/A}$ is the two-term complex
$[T_Y\to f^*T_A]$ in degrees $0,1$, so
$T^2_{Y/A}(\mathcal O_Y)=H^1(Y,N_f)$.
Calabi--Vesentini and the cokernel sequence give
$H^1(\mathcal O_Y)^{\oplus g}\hookrightarrow H^1(N_f)$.
The target specified in Buchweitz--Flenner, Proposition~7.21,
Lemma~7.22 and the definition before Theorem~7.23, is
\[
 \bigoplus_{r=0}^{g-2}H^{r+2}(A,\Omega_A^r).
\]
Its dimension is
$\sum_{r=0}^{g-2}\binom gr\binom g{r+2}=\binom{2g}{g-2}$
by Vandermonde. Rank--nullity proves the assertion.
\end{proof}

\begin{corollary}[S-06: Failure at sufficiently deep allowed levels]
\label{cor:semireg-level-failure}
Fix a nonzero 032 surface construction $f:Y\to A$. There are
connected finite \'etale covers $\pi_j:Y_j\to Y$, preserving its
four pulled-back Hodge maps and nonzero period, for which
$\tau_{f\pi_j}$ is not injective. Such covers may be chosen among
a cofinal system of principal congruence levels.
\end{corollary}
\begin{proof}
The nonzero period implies $H^1(Y,\mathbf C)\ne0$.
Venkataramana, \emph{Virtual Betti numbers of compact locally
symmetric spaces}, Theorem~1, proves that finite-dimensional
direct-limit cohomology over finite covers equals the cohomology
of the compact dual. Here the compact dual is $\mathbf P^2$,
whose first cohomology is zero. Pullback under finite covers is
injective over $\mathbf Q$ by transfer, so the direct limit is
nonzero and therefore infinite-dimensional. Hence $q(Y_j)$ is
unbounded, and the preceding proposition applies with fixed $A$.
The congruence version is explicitly identified in the final Remark
on p.~3 of that paper. Its proof uses the same rational-conjugation
action, since rational conjugation and finite intersection preserve
congruence subgroups; principal subgroups are cofinal.
Persistence of the period follows from the first proposition.
\end{proof}

This is a statement about the deformation problem of a morphism.
The factorization
$\tau_f=\sigma_{Rf_*\mathcal O_Y}\circ\psi_f$
does not imply that a kernel of $\tau_f$ survives in the perfect
complex obstruction space: $\psi_f$ itself can have a kernel.
Nor does the corollary settle Bloch semiregularity of a particular
smooth reduced image.

\begin{proposition}[S-07: Descent of the cover obstruction space]
\label{prop:semireg-cover-invariants}
For a finite \'etale Galois cover $\pi:Y'\to Y$ with deck group $H$,
$N_{f\pi}=\pi^*N_f$ and
$H^1(Y',N_{f\pi})^H=H^1(Y,N_f)$.
\end{proposition}
\begin{proof}
Flat pullback identifies the tangent cokernel sequences.
Finite pushforward and the projection formula identify
$\pi_*\pi^*N_f=N_f\otimes\pi_*\mathcal O_{Y'}$.
In characteristic zero averaging over $H$ splits off precisely
the invariant summand $N_f$. Cohomology commutes with this direct
summand, giving the assertion.
\end{proof}

\subsection{Derived replacements and their remaining gate}

\begin{proposition}[S-09: An integral Chern-component projector]
\label{prop:semireg-k-projector}
Let $A$ be an abelian $g$-fold and $\xi\in K_0(A)$ with
$\operatorname{ch}_j(\xi)=0$ for $j<p$. For $m\geq2$ put
\[
 M=[m]^*,\qquad
 D=\prod_{j=p+1}^g(m^{2j}-m^{2p}),\qquad
 \eta=\prod_{j=p+1}^g(m^{2j}\operatorname{id}-M)\xi.
\]
Then $\eta\in K_0(A)$ and
$\operatorname{ch}(\eta)=D\operatorname{ch}_p(\xi)$ in rational
Betti cohomology.
\end{proposition}
\begin{proof}
The operator $M$ multiplies $H^{2j}(A,\mathbf Q)$ by $m^{2j}$.
Thus every higher Chern component is annihilated by its corresponding
factor, lower components were already zero, and the component of
index $p$ is multiplied by $D$. The polynomial has integer
coefficients. Integral K-classes are represented by perfect complexes
on the smooth projective variety, using finite locally free
resolutions and shifts to implement negative coefficients.
\end{proof}

Starting from an integral codimension-$p$ cycle of class $N_0t$,
one may use $\xi=\sum a_i[\mathcal O_{Z_i}]$; its leading Chern
character is the fundamental cycle. Alternatively, resolve each
$Z_i$ and use the proper pushforward of its structure sheaf;
Grothendieck--Riemann--Roch gives the same leading component.
Adding line-bundle classes gives the useful character
$Nt+\sum_a b_a e^{a\theta}$, which remains Hodge on the polarized
Route-2 family. This constructs a K-class, not a semiregular object.

\begin{proposition}[S-08: Multiplicity creates trace kernels]
\label{prop:semireg-multiplicity}
If $E$ is perfect and $\operatorname{Ext}^2(E,E)\ne0$, then
$E^{\oplus m}$ is not fully semiregular for any $m\geq2$.
\end{proposition}
\begin{proof}
Its degree-two self-Ext group is
$\operatorname{Ext}^2(E,E)\otimes\operatorname{End}(\mathbf C^m)$.
The Atiyah class is $\operatorname{at}(E)\otimes\operatorname{id}$;
therefore every semiregularity component is the corresponding
component for $E$ tensored with the matrix trace.
A nonzero degree-two Ext class tensored with a nonzero traceless
matrix lies in the kernel.
\end{proof}

\begin{proposition}[S-10: Full Fourier--Mukai transport]
\label{prop:semireg-full-fm}
For the fixed enhanced Fourier--Mukai equivalence in the supplied
032 companion \S5.2, the full trace maps satisfy
\[
 \sigma_{\Phi Q}\circ\Phi_{\rm Ext}
   =\Phi^H_{HH}\circ\sigma_Q
 \quad\hbox{on }\operatorname{Ext}^2_A(Q,Q).
\]
Thus full semiregularity is invariant under this equivalence.
\end{proposition}
\begin{proof}
Use the companion's enhanced adjunctions and notation.
Put
$b_Q(v)=\operatorname{Tr}_{Q,H_A[2]}(A_Q^H[2]\circ v)
\in HH_{-2}(A)$.
Equations (5.20), (5.24) identify its HKR image with $\sigma_Q(v)$.
For $\delta\in T^{-2}(A)$, trace-compatible adjunction gives
\[
 \langle b_Q(v),\delta\rangle
   =\operatorname{Tr}_{A,Q}(\delta_Q[2]\circ v).
\]
Indeed, if $\delta':H_A\to\omega_A[g-2]$ is the adjoint coefficient
map, move $\delta'[2]$ through the partial trace by (5.24);
evaluation then gives the right side. This is exactly (5.37)
with the displayed even shift.
The convolution proof of (5.33), with the same inverse adjunction
data and shifted by two, gives
$(\Psi^{-2}\delta)_{\Phi Q}=\theta_Q[-2]\Phi(\delta_Q)$.
By the trace preservation in (5.31),
\[
 \langle b_{\Phi Q}(\Phi_{\rm Ext}v),\Psi^{-2}\delta\rangle
  =\operatorname{Tr}_{A,Q}(\delta_Q[2]\circ v)
  =\langle b_Q(v),\delta\rangle.
\]
Definition (5.34) and nondegeneracy of the Serre pairing imply
$b_{\Phi Q}(\Phi_{\rm Ext}v)=\Phi_*^{HH}b_Q(v)$.
Apply (5.35) and the fixed HKR maps. The two vertical maps
are isomorphisms, proving the assertion about kernels.
This proves a refinement; Proposition~5.2 itself only asserted
transport on the obstruction-evaluation image.
\end{proof}

\begin{proposition}[S-11: An ample complete intersection need not be semiregular]
\label{prop:semireg-ci-counterexample}
Let $A$ be an abelian fourfold, $L$ very ample, and
$Y=D_1\cap D_2$ a smooth transverse intersection of two members
of $|L|$. Then $h^1(N_{Y/A})=24$, whereas
$h^3(A,\Omega_A^1)=16$. Thus its Bloch map has a kernel of
dimension at least eight.
\end{proposition}
\begin{proof}
The twisted Koszul sequence is
$0\to L^{-1}\to\mathcal O_A^{\oplus2}\to\mathcal I_Y\otimes L\to0$.
Kodaira vanishing and $K_A=\mathcal O_A$ give $H^j(A,L)=0$ for
$j>0$; Serre duality gives $H^j(A,L^{-1})=0$ for $j<4$.
Thus
$H^2(A,\mathcal I_Y\otimes L)=H^2(A,\mathcal O_A)^{\oplus2}$.
The restriction sequence gives
$H^1(Y,L|_Y)=H^2(A,\mathcal I_Y\otimes L)$, of dimension twelve.
Transversality identifies $N_{Y/A}=(L|_Y)^{\oplus2}$.
The Bloch target has dimension
$\binom43\binom41=16$, so rank--nullity finishes the proof.
\end{proof}

\begin{proposition}[S-12: A pure Weil character obstruction for coherent sheaves]
\label{prop:semireg-pure-weil-no-go}
Let $A_0$ be a projective member of Voisin's four-dimensional
marked $\mathbf Q(i)$-Weil torus family, and let $t\ne0$ be a
rational Weil class. No fully Buchweitz--Flenner-semiregular
coherent sheaf $F$ on $A_0$ satisfies
\[
 \operatorname{ch}_1(F)=\operatorname{ch}_3(F)=0,\qquad
 \operatorname{ch}_2(F)=t.
\]
The rank and top-degree character are unrestricted.
\end{proposition}
\begin{proof}
The unpolarized Weil torus parameter space is smooth near $A_0$;
Voisin, \emph{A counterexample to the Hodge conjecture extended
to K\"ahler varieties}, \S3, equations (3.1)--(3.8), describes it.
Lemma~3.1 shows that the rational Weil space remains of Hodge
type $(2,2)$. Hence the components with indices $I=\{1,2,3\}$
in the assertion remain horizontal Hodge classes. In dimension
four these are precisely the components tested by full
semiregularity.
Buchweitz--Flenner Theorem~5.1 and its proof on p.~179 would
therefore extend $F$ as a coherent flat sheaf on a neighborhood.
Voisin Proposition~3.2 and Proposition~2.1 show that very general
fibres in that neighborhood have $\operatorname{NS}=0$ and
$c_2(G)=0$ for every coherent analytic sheaf $G$.
On such a fibre $c_1(F)=c_2(F)=0$, whence
$\operatorname{ch}_2(F)=0$, contrary to the nonzero horizontal
class $t$.
\end{proof}

This last proposition is a coherent-sheaf statement for the
specified fourfold family. Extension to arbitrary perfect complexes
would require an analytic effectivity theorem for their formal
deformations; algebraic/formal Pridham alone is not that theorem.
Allowing a nonzero polarization term in the Chern character
avoids this particular contradiction.

\begin{proposition}[S-13: One exponential polarization term does not repair the coherent gate]
\label{prop:semireg-single-exponential-no-go}
In the fourfold setting of S-12, let $L$ be a polarization whose
Rosati involution restricts to conjugation on $\mathbf Q(i)$, and
write $\theta=c_1(L)$. If $a\in\mathbf Z$, $\lambda\in\mathbf Q^\times$,
and $t\ne0$ is a rational Weil class, no fully semiregular coherent
sheaf $F$ has
\[
 \operatorname{ch}(F)=r e^{a\theta}+\lambda t+c[\mathrm{pt}]
\]
for any $r,c\in\mathbf Q$.
\end{proposition}
\begin{proof}
Write $H^1(A_0,\mathbf C)=V_\sigma\oplus V_{\bar\sigma}$.
The Rosati condition places $\theta$ in
$V_\sigma\otimes V_{\bar\sigma}$, whereas the complexified Weil
space is
$\bigwedge^4V_\sigma\oplus\bigwedge^4V_{\bar\sigma}$.
Both eigenspaces have dimension four. Consequently
$\theta\smile t=0$, and hence $e^{-a\theta}t=t$.
Tensoring by $L^{-a}$ is a Fourier--Mukai autoequivalence.
The functoriality in Perry Remarks~2.3 and~2.5, p.~5, preserves
full semiregularity under this equivalence, as in S-10.
It preserves coherence as well.
For $F'=F\otimes L^{-a}$ the Chern-character formula gives
\[
 \operatorname{ch}(F')=r+\lambda t+c[\mathrm{pt}].
\]
Thus $\operatorname{ch}_1(F')=\operatorname{ch}_3(F')=0$ and
$\operatorname{ch}_2(F')=\lambda t\ne0$, contradicting S-12.
\end{proof}

The integer condition on $a$ ensures that this normalization is an
actual line-bundle twist. The conclusion does not exclude all
polarization polynomials: it excludes the displayed single
exponential, including an arbitrary additional top-degree term.

\begin{proposition}[S-14: An equivariant perfect-complex gate implies bounded division]
\label{prop:semireg-perry-bounded-division}
Let $\pi:A\to S$ be an abelian scheme over a smooth irreducible
complex algebraic variety of finite type, with a relatively very
ample line bundle $H$ and $\theta=c_1(H)$. Let
$t\in H^0(S,R^{2p}\pi_*\mathbf Q(p))$ be a flat section whose values
are Hodge classes, where $0\leq p\leq\dim(A/S)$.
Suppose that for one $s_0\in S(\mathbf C)$ there exist:
\begin{enumerate}
\item a finite subgroup $G\subset A_{s_0}\times A_{s_0}^{\vee}$
with a specified coherent action on $\operatorname{Perf}(A_{s_0})$
by translations and tensoring by line bundles;
\item a weakly $G$-semiregular perfect complex $E_0$ with
$\operatorname{Ext}^{i}(E_0,E_0)=0$ for all $i<0$;
\item an algebraic class $B_0\in H^2(A_{s_0},\mathbf Q(1))$,
a nonzero $\lambda\in\mathbf Q$, and $P\in\mathbf Q[x]$ such that
\[
 e^{B_0}\operatorname{ch}(E_0)=\lambda t_{s_0}+P(\theta_{s_0})
\]
in the full even rational cohomology.
\end{enumerate}
Then there is an integer $C\geq1$ such that, for every
$s\in S(\mathbf C)$, there are an integer $1\leq N_s\leq C$
and effective integral codimension-$p$ cycles $Z_s^+,Z_s^-$,
each of $H$-degree at most $C$, with
\[
 N_s t_s=\operatorname{cl}(Z_s^+)-\operatorname{cl}(Z_s^-).
\]
In particular, there are single integers $N,D\geq1$ for the entire
family such that $Nt_s$ admits such representatives of degrees
at most $D$ on every fibre.
A sufficient form of the second hypothesis is a specified
$G$-equivariant structure on $E_0$ for which the ordinary full
semiregularity map is injective on
$\operatorname{Ext}^2(E_0,E_0)^G$, together with the same vanishing
of negative self-Ext groups.
\end{proposition}
\begin{proof}
The class
$w=\lambda t+P(\theta)$ is a global flat section of the full even
cohomology local system, and every component has the required
Hodge type. It extends
$w_0=e^{B_0}\operatorname{ch}(E_0)$.
Apply Alexander Perry, \emph{The semiregularity theorem for
equivariant noncommutative varieties},
arXiv:2604.00511v2 (24 June 2026), Theorem~1.2, p.~3.
Its second conclusion states that $w_s$ is algebraic for every
$s\in S(\mathbf C)$. Since $P(\theta_s)$ is algebraic and
$\lambda\ne0$, each $t_s$ is algebraic.
DIV1(2), applied to the fixed smooth projective family and $H$,
now supplies a single bound $C$ for the signed integral complexity.
This is exactly the displayed conclusion.
For a common multiplier take $N=\operatorname{lcm}(1,\ldots,C)$
and multiply the cycles on fibre $s$ by $N/N_s$.
Their degrees are at most $D=NC$.
Finally, Perry Definition~2.6 and Lemma~2.7, pp.~5--6, show that
the stated invariant-injectivity criterion implies weak
$G$-semiregularity over $\mathbf C$.
\end{proof}

The gate requires the entire displayed Chern character to remain
Hodge, a coherent categorical action, and vanishing of negative
self-Ext groups. A virtual K-class alone supplies none of these
deformation properties. The theorem supplies algebraicity on
every fibre; no properness assertion for a moduli space of
perfect complexes is used. The bound obtained from DIV1 is
existential. Here $G$ is a subgroup of the target's translation
and degree-zero line-bundle group, not an arbitrary deck group
of the source surface.

\begin{proposition}[S-15: Canonical source-cover equivariance retains a full kernel]
\label{prop:semireg-cover-equivariance}
Let $f:Y\to A$ be a proper morphism from a smooth projective
variety to an abelian variety, and put $E=Rf_*\mathcal O_Y$.
\begin{enumerate}
\item If $\rho:Y'\to Y$ is a finite \'etale Galois cover with
group $K$, the canonical $K$-equivariant lift of
$E'=R(f\rho)_*\mathcal O_{Y'}$ for the trivial action of $K$
on $\operatorname{Perf}(A)$ can be semiregular only if $E$
is fully semiregular.
\item Suppose $\rho$ is obtained by base change from an isogeny
$q:A'\to A$. Set
$G=\ker(q^*:A^\vee\to(A')^\vee)$, acting on
$\operatorname{Perf}(A)$ by the coherently chosen tensor
line bundles in $q_*\mathcal O_{A'}$.
The canonical $G$-equivariant lift of
$q_*q^*E=R(f\rho)_*\mathcal O_{Y'}$ can be semiregular only if
$E$ is fully semiregular.
\end{enumerate}
These assertions concern the canonical lifts supplied by the
covers; they do not exclude every possible new equivariant
structure on the underlying objects.
\end{proposition}
\begin{proof}
For (1), finite pushforward and characteristic-zero averaging give
\[
 (E')^K
 =Rf_*(\rho_*\mathcal O_{Y'})^K
 =Rf_*\mathcal O_Y=E.
\]
Because the categorical action is trivial,
$\operatorname{Perf}(A)^K$ is the finite product of copies of
$\operatorname{Perf}(A)$ indexed by irreducible representations
of $K$. The factor belonging to the trivial representation of
the canonical lift is $E$. Self-Ext groups, Hochschild homology,
and the semiregularity map decompose by these factors.
Semiregularity of the lift therefore implies that of $E$.

For (2), let $R=q_*\mathcal O_{A'}$. The character decomposition
of the finite abelian torsor gives
$R=\bigoplus_{P\in G}P$, with multiplication isomorphisms
$P\otimes Q\cong PQ$. These provide the indicated coherent
tensor action. A $G$-equivariant object is equivalently an
$R$-module object; since $q$ is finite \'etale this identifies
$\operatorname{Perf}(A)^G$ with $\operatorname{Perf}(A')$.
Under this equivalence the canonical lift of $q_*q^*E$
corresponds to $q^*E$.
Flat base change for the square defining $Y'$ identifies the
underlying object with the stated source-cover pushforward.

The pullback map
\[
 q^*:\operatorname{Ext}^2_A(E,E)
 \longrightarrow
 \operatorname{Ext}^2_{A'}(q^*E,q^*E)
\]
is injective: by adjunction its target is
$\operatorname{Ext}^2_A(E,E\otimes R)$, and the unit
$\mathcal O_A\to R$ has the retraction
$(\deg q)^{-1}\operatorname{Tr}_q$.
Semiregularity is natural for the enhanced functor $q^*$,
by Perry Remark~2.3, p.~5. Thus pullback carries
$\ker\sigma_E$ injectively into $\ker\sigma_{q^*E}$.
Semiregularity of the canonical lift would make the latter
kernel zero, proving the assertion.
\end{proof}

For an ordinary source deck transformation $k$ one has
$f\rho k=f\rho$. Consequently an induced translation of the
target satisfying this equivariance must be the identity.
An ordinary deck action therefore supplies internal
automorphisms of the pushed-forward object; it does not by
itself specify the nontrivial target action required in S-14.
Part (2) treats a special case where target tensor symmetries
really do arise, and shows why their canonical lift still
does not remove a pre-existing full semiregularity kernel.

There is also an independent primary-source confirmation of S-10:
Perry Remarks~2.3 and~2.5, p.~5, give functoriality of the
semiregularity morphism on the entire mapping complex and identify
its degree-two part with Buchweitz--Flenner under HKR.
An equivalence induces isomorphisms on both mapping complexes and
Hochschild homology, so the full kernel is invariant. This argument
does not depend on a raw versus Todd-modified HKR convention;
for abelian varieties the Todd class is one.

\begin{proposition}[S-16: A polarization exponential realizes a K-class by a vector bundle]
\label{prop:semireg-kclass-bundle-realization}
Let $X$ be a smooth projective variety, let $L$ be ample, and let
$\xi\in K_0(X)$. There exist integers $m,r\geq1$ and a vector
bundle $V$ such that
\[
 [V]=\xi+r[L^{-m}]\quad\hbox{in }K_0(X).
\]
In particular, if $\operatorname{ch}(\xi)=\lambda t$, then
$\operatorname{ch}(V)=r e^{-m c_1(L)}+\lambda t$.
In the setting of S-13, with $\lambda t\ne0$, every vector bundle
with this character fails full semiregularity.
\end{proposition}
\begin{proof}
On a smooth projective variety a K-class is a difference
$\xi=[U]-[W]$ of vector bundles: take finite locally free
resolutions and group their even and odd terms.
For sufficiently large $m$, the vector bundle $W\otimes L^m$
is generated by finitely many global sections.
The precise ampleness criterion is Stacks Project,
Proposition~28.27.13(7)--(8), tag 01Q3.
Choose a surjection
\[
 0\longrightarrow K\longrightarrow\mathcal O_X^{\oplus r}
   \longrightarrow W\otimes L^m\longrightarrow0 .
\]
One can add a redundant zero section to ensure $r\geq1$.
The kernel $K$ is locally free, since the surjection onto a
vector bundle splits locally. After tensoring by $L^{-m}$,
the K-group relation is
$[K\otimes L^{-m}]=r[L^{-m}]-[W]$.
Thus $V=U\oplus(K\otimes L^{-m})$ has the required class.
The Chern-character formula follows by additivity, and the
last assertion is S-13 with $a=-m$ and $c=0$.
\end{proof}

This gives an actual coherent vector-bundle representative
after one exponential addition. It supplies neither full nor
equivariant semiregularity; in the indicated fourfold case the
full version is ruled out by S-13.

\begin{proposition}[S-17: Tensor--determinant untwisting after a positive-rank deformation]
\label{prop:semireg-tensor-determinant}
Let $\pi:\mathcal X\to X$ be a $\mu_M$-gerbe over a connected
smooth complex scheme, and let $F$ be a weight-one perfect complex
of positive virtual rank $r$. Then
\[
 F^{\otimes^{\mathbf L}r}\otimes\det(F)^{-1}
 \simeq\pi^*K
\]
for an ordinary perfect complex $K$ on $X$.

In the family setting and under all the hypotheses of S-14,
suppose $\dim A_s=2n$, $n\ge2$, that $t$ is a nonzero relative
Weil class of codimension $n$, and that
\[
 w_0=e^{B_0}\operatorname{ch}(E_0)
     =P(\theta_0)+\lambda t_0,
 \qquad P(0)=r>0,\quad\lambda\ne0,\quad\theta_0t_0=0.
\]
Put $p_1=[x]P(x)$. On the local deformation neighborhood supplied
by Perry, the descended complex has character, through
codimension $n$,
\[
 \operatorname{ch}(K_s)
 \equiv e^{-p_1\theta_s}P(\theta_s)^r+r^r\lambda t_s.
 \tag{S17.1}
\]
There are fixed integers $R,m\ge1$ and $a,k\ge0$ from a finite
bundle presentation and generation twists for which the following
numbers give a uniform signed-cycle bound on the entire base.
Set
\[
 Q(x)=e^{kx}\bigl(e^{-p_1x}P(x)^r+R e^{-mx}+a\bigr),
 \qquad\mu=r^r\lambda,
\]
\[
 C_n=[u^n]\exp\!\left(\sum_{j=1}^n
     (-1)^{j-1}(j-1)![x^j]Q(x)u^j\right),
 \qquad\beta=(-1)^{n-1}(n-1)!\mu .
\]
If $d\ge1$ clears the denominators of $C_n$ and $\beta$, and
$v=\int_{A_s}\theta_s^{2n}$, then
\[
 N=|d\beta|>0,\qquad D=dC_nv>0
 \tag{S17.2}
\]
are fixed integers such that $Nt_s$ is the difference of two
effective integral cycles, each of polarization degree at most
$D$, on every fibre. A negative nonzero virtual rank is reduced
to this case by shifting the initial and deformed complexes.
\end{proposition}
\begin{proof}
The determinant functor for perfect complexes exists and commutes
with base change by Knudsen--Mumford, \emph{The projectivity of the
moduli space of stable curves I}, Theorem~2, p.~42, with the
extended determinant definition on pp.~40--42. Apply it on an
\'etale atlas of the gerbe. Functoriality and base-change
compatibility glue the local determinants to a line bundle
$\det(F)$ on $\mathcal X$. Locally a finite locally free
presentation gives
\[
 \det(F)=\bigotimes_i\det(F^i)^{(-1)^i}.
\]
The inertia element $\zeta\in\mu_M$ acts on this line by
$\zeta^{\sum_i(-1)^i\operatorname{rk}F^i}=\zeta^r$.
Consequently $F^{\otimes^{\mathbf L}r}\otimes\det(F)^{-1}$
has weight zero. Weight-zero quasi-coherent complexes on a
$\mu_M$-gerbe are precisely pullbacks from the base: after an
\'etale trivialization this is the weight-zero summand of the
representation category of $\mu_M$, and the identifications
descend. Perfectness is local for this topology, proving the
first assertion.

Use the weight-one deformation $F$ on the gerbe in Perry's proof
of Theorem~1.2, \S6.2, pp.~26--28. At the central fibre its
normalization has character $w_0$ and first component
$w_{0,1}=p_1\theta_0$. Directly, cancellation of the $B$-field
in the tensor--determinant expression gives
\[
 \operatorname{ch}(K_0)
 =\operatorname{ch}(E_0)^r
       e^{-c_1(E_0)}
 =w_0^r e^{-w_{0,1}} .
\]
This also follows from the central root description of the gerbe:
the root-line factors cancel between the $r$ tensor factors and
the determinant. Since $n\ge2$, $w_{0,1}=p_1\theta_0$ contains
no Weil term. Through codimension $n$, the only occurrence of
$t_0$ in $w_0^r$ is
$rP(0)^{r-1}\lambda t_0=r^r\lambda t_0$.
All positive powers of $\theta_0$ annihilate $t_0$.
This proves (S17.1) at the central fibre; flatness of Chern classes
proves it over the connected deformation neighborhood.
Terms in higher codimension, including a possible $t_0^2$ term,
are not used in this formula or in the computation of $c_n$.

Shrink the neighborhood to a smooth affine base $U$.
Its total space $A_U$ is smooth and quasi-projective, and its
relatively very ample bundle $H$ is ample over this affine base.
Stacks Project Lemma~36.37.2, tag 0F8F, represents a perfect
complex by a bounded complex of vector bundles on a scheme
with the resolution property. Thus write
$[K]=[V^+]-[V^-]$. For sufficiently large $m$, choose a finite
surjection $\mathcal O^{\oplus R}\to V^-\otimes H^m$.
Its kernel $J$ is locally free, and
\[
 [V^+\oplus(J\otimes H^{-m})]=[K]+R[H^{-m}].
\]
Adding redundant sections ensures $R\ge1$.
This is the relative version of S-16, using the finite global
generation criterion of Stacks Proposition~28.27.13(7)--(8),
tag 01Q3. Add $a$ trivial summands to make the rank at least $n$,
and twist by a sufficiently large $H^k$ to obtain a globally
generated vector bundle $V$. Its character through codimension
$n$ is $Q(\theta_s)+\mu t_s$.

The identity
\[
 c_u(V)=\exp\!\left(\sum_{j\ge1}
             (-1)^{j-1}(j-1)!\operatorname{ch}_j(V)u^j\right)
\]
gives $c_n(V_s)=C_n\theta_s^n+\beta t_s$.
If $q=\operatorname{rk}V$, the degeneracy cycle of $q-n+1$
general sections is effective of codimension $n$, or empty,
and represents $c_n(V_s)$; see Fulton, \emph{Intersection
Theory}, Theorem~14.4 and Examples~14.4.2--3, pp.~254--256.
Call it $Z_s$. Its Weil component $\beta t_s$ is nonzero,
so the cycle is nonempty. Since $\theta_st_s=0$, its degree is
$C_nv$, which is positive. Let $W_s$ be a complete intersection
of $n$ general members of $|H_s|$; it represents $\theta_s^n$
and has degree $v$. Then
\[
 d\beta t_s=\operatorname{cl}(dZ_s)
                   -\operatorname{cl}(dC_nW_s).
\]
Both cycles are integral and effective of degree $D$.
If $\beta<0$, interchange them. This proves (S17.2) over a
nonempty \'etale image, which is a dense Zariski open subset
of the irreducible base. The fixed-$N$, fixed-$D$ locus of such
signed representatives is Zariski closed by the proper
parameter-space argument in DIV1. It therefore equals the
entire base, preserving these same constants.

Finally, a shift changes the signs of the virtual rank and
Chern character, preserves the self-Ext groups and
semiregularity, and shifts the existing deformation.
This reduces negative rank to positive rank without adding
a new existence hypothesis.
\end{proof}

This construction is used only after the successful twisted
deformation exists. No semiregularity of $K$ or of $V$ is claimed.
The Adams construction DER4 removes the nonzero-rank restriction;
S-17 is a separate, simpler tensor--determinant alternative.

\begin{proposition}[S-18: Canonical induction preserves a full semiregularity kernel]
\label{prop:semireg-induction-kernel}
Let $\mathcal C$ be a smooth proper complex linear enhanced
category with a coherent action of a finite group $G$ by exact
autoequivalences, and assume its equivariant category
$\mathcal C^G$ is smooth proper. Let
\[
 \operatorname{Ind}:\mathcal C\longrightarrow\mathcal C^G
\]
be induction, whose underlying object is
$\bigoplus_{g\in G}gE$, with the canonical permutation
equivariant structure. Induction is injective on every morphism
space. In particular it induces an injection
\[
 \ker\sigma_E\hookrightarrow
       \ker\sigma_{\operatorname{Ind}(E)}.
\]
Thus a nonzero full semiregularity kernel cannot be removed by
canonical induction, including for actions combining
translations and tensoring by degree-zero line bundles on an
abelian variety.
\end{proposition}
\begin{proof}
For $u:E\to E'[j]$, the underlying map of
$\operatorname{Ind}(u)$ is the block diagonal map
$\bigoplus_{g\in G}g(u)$.
If it is zero, compose with the inclusion of the identity
summand of the source and the projection to the identity
summand of the target. The coherence isomorphism for the
identity element identifies this composite with $u$, so $u=0$.
In particular, if $\operatorname{Ind}(u)$ is zero in the
equivariant category, forgetting equivariance gives zero and
the same argument applies. This proves injectivity in every
degree, without requiring the forgetful functor to be full.

Induction is an exact enhanced functor, so naturality of the
full semiregularity morphism gives
\[
 \sigma_{\operatorname{Ind}(E)}(\operatorname{Ind}(u))
       =\operatorname{Ind}_*\sigma_E(u).
\]
This is Perry Remark~2.3, p.~5; Remark~2.5 identifies the
geometric map with Buchweitz--Flenner under HKR.
Consequently a nonzero element of $\ker\sigma_E$ has a
nonzero image lying in the displayed target kernel.
\end{proof}

S-18 rules out semiregularity of the canonical induced lift.
It does not rule out a different equivariant structure on the
same underlying object: weak equivariant semiregularity is
an existence condition over possible lifts, as specified in
Perry Definition~2.6, pp.~5--6.

The qualitative fourfold secant obstruction and its free-translation
variant are prior results of rStation paper~37, version~4,
Theorems~A--C and~$B_G$. The next two applications add
explicit bounds for the supplied Markman constructions.

\begin{proposition}[S-19: The isolated-point genus-four secant seed fails at every parameter]
\label{prop:semireg-markman-genus-four-points}
Consider the first genus-four construction of Markman
Example~8.2.3 in the supplied 31 May 2025 source. Thus $X$ is
a genus-four Jacobian with principal polarization $\theta$,
$d\ge1$, and $G\subset X$ is a cyclic translation group of
order $\ell=d+1$. The $G$-invariant reduced subscheme $Z$ is
the union of $\ell$ translates of $W_2$ in the stated general
position and $2d\ell$ isolated reduced points away from them.
Let $q:X\to Y=X/G$, let $\bar Z=Z/G$, and set
$E=\mathcal I_{\bar Z}$. Then
\[
 \chi(E,E)=8d,\qquad
 \operatorname{ext}^1(E,E)\ge8d+4,\qquad
 \operatorname{ext}^2(E,E)\ge24d+6>28.
 \tag{S19.1}
\]
In particular $E$ is not fully semiregular. The original ideal
$\mathcal I_Z$ is not weakly $G$-semiregular for this translation
action, for any choice of linearization. The same holds for
$\mathcal I_Z(\theta)$ under the action conjugated by tensoring
with $\mathcal O_X(\theta)$.
\end{proposition}
\begin{proof}
The quotient is finite \'etale of degree $\ell$. Its isolated
point set has $2d$ reduced points, and
$q^*\mathcal I_{\bar Z}=\mathcal I_Z$.
Markman's character calculation in Example~8.2.3 gives
\[
 w:=\operatorname{ch}(\mathcal I_Z(\theta))
   =1+\theta-\frac d2\theta^2-\frac d6\theta^3
          +\frac{d^2}{24}\theta^4 .
\]
Since $\int_X\theta^4=24$, direct multiplication yields
$\int_X w^\vee w=8d(d+1)$. Tensoring by the theta line bundle
does not change this Euler pairing. Riemann--Roch and the
\'etale degree formula therefore give $\chi(E,E)=8d$.

We verify the asserted tangent-space lower bound. On a smooth
variety, the endomorphism sheaf of an ideal of codimension at
least two is $\mathcal O$: an endomorphism is multiplication
by a regular function away from that locus, and normality
extends the function. The unit and perfect trace give a split
inclusion
\[
 \mathcal O_Y\longrightarrow R\mathcal Hom(E,E)
       \longrightarrow\mathcal O_Y
\]
whose composite is the virtual rank, namely one. Its complement
has no degree-zero cohomology and no negative cohomology.
Consequently
\[
 \operatorname{Ext}^1(E,E)
 \cong H^1(Y,\mathcal O_Y)
       \oplus H^0(Y,\mathcal Ext^1(E,E)).
\]
At each isolated reduced point the ideal stalk is the maximal
ideal $\mathfrak m$ of a regular local ring $R$ of dimension
four. Apply $\operatorname{Hom}_R(\mathfrak m,-)$ to
$0\to\mathfrak m\to R\to\mathbf C\to0$.
Here $\operatorname{Hom}_R(\mathfrak m,R)=R$ maps to zero in
$\operatorname{Hom}_R(\mathfrak m,\mathbf C)$, and
\[
 \operatorname{Ext}^1_R(\mathfrak m,R)
   =\operatorname{Ext}^2_R(\mathbf C,R)=0.
\]
Thus
$\operatorname{Ext}^1_R(\mathfrak m,\mathfrak m)
\cong(\mathfrak m/\mathfrak m^2)^\vee$ has dimension four.
The isolated supports contribute independent skyscraper
summands. Since $h^1(\mathcal O_Y)=4$, this proves
$\operatorname{ext}^1(E,E)\ge4+8d$.

The sheaf $E$ is simple. Serre duality on the abelian fourfold
gives $\operatorname{ext}^4=\operatorname{ext}^0=1$ and
$\operatorname{ext}^3=\operatorname{ext}^1$; there are no
self-Ext groups outside degrees zero through four. Hence
\[
 \operatorname{ext}^2(E,E)
 =\chi(E,E)-2+2\operatorname{ext}^1(E,E)
 \ge24d+6.
\]
The full semiregularity target on an abelian fourfold has
dimension
\[
 h^2(\mathcal O_Y)+h^3(\Omega_Y^1)+h^4(\Omega_Y^2)
 =6+16+6=28,
\]
proving failure of injectivity.

Equivariant descent identifies
$\operatorname{Perf}(X)^G$ with $\operatorname{Perf}(Y)$
and the canonical lift of $\mathcal I_Z$ with $E$.
Moreover $\operatorname{Aut}(\mathcal I_Z)=\mathbf C^*$,
so any two $G$-linearizations differ by a character of $G$.
Their descents differ by tensoring with the associated
degree-zero torsion line bundle on $Y$. Such a tensor
equivalence preserves the full kernel by S-10.
Thus no linearization is semiregular in the invariant category.
Conjugating this action by the theta tensor equivalence gives
the last assertion; the conjugated action consists of target
translations and tensoring by degree-zero line bundles.
\end{proof}

\begin{proposition}[S-20: The fourfold adjunction test excludes the partial-normalization seed]
\label{prop:semireg-fourfold-partial-normalization}
On an abelian fourfold $Y$, let $E$ be perfect and let
$c_E^2:HT^2(Y)\to H\Omega^{-2}(Y)$ be contraction with
$\operatorname{ch}(E)$, in the conventions of the supplied
032 companion, \S5.1. If $E$ is fully semiregular, then
\[
 \operatorname{ext}^2(E,E)=\operatorname{rank}c_E^2.
 \tag{S20.1}
\]
For the cyclic partial-normalization construction in Markman
Example~8.2.3, let $\bar E$ be the descended complex
\[
 [\mathcal O_Y\longrightarrow
       \bar\nu_*\mathcal O_{\widetilde Z/G}]
\]
in degrees zero and one, where four of the six transverse
meeting points for each pair of translated $W_2$ surfaces
are normalized, and the collection of choices is $G$-invariant.
Then, for every $d\ge1$,
\[
 \operatorname{rank}c_{\bar E}^2=12,
 \qquad
 \operatorname{ext}^2(\bar E,\bar E)\ge8d+6>12.
 \tag{S20.2}
\]
Thus this complex is not fully semiregular, and its pullback
is not weakly $G$-semiregular for the specified free translation
action, with any linearization. In fact
\[
 \dim\ker\sigma_{\bar E}\ge4d-3.
 \tag{S20.3}
\]
\end{proposition}
\begin{proof}
The adjunction criterion is the previously known CY4 test
(rStation paper~37, version~4, Theorem~A). In the conventions
used here, 032 equations (5.20)--(5.24) identify $\sigma_E$
with the Serre transpose of $\operatorname{ob}_E$, and
Proposition~5.1, equation (5.9), gives
$\sigma_E\operatorname{ob}_E=c_E^2$. The Todd correction is
one on an abelian variety. Thus injectivity of $\sigma_E$
forces surjectivity of $\operatorname{ob}_E$ and proves
(S20.1).

The imaginary-quadratic rank calculation is likewise the
one used in that prior criterion. For the character
$w=\alpha+\beta$ of S-19, put $U=H^1(\mathcal O_Y)$,
$A=M\mathbin{\lrcorner}\theta$,
$B=\tfrac12\nu\mathbin{\lrcorner}\theta^2$, and
$D=\eta-dB$, where
$(\eta,M,\nu)\in\wedge^2U\oplus(U\otimes H^0(T_Y))
\oplus\wedge^2H^0(T_Y)$.
The first two output blocks are
$D+A$ and $\theta\wedge(D-dA)$.
Their coefficient determinant is $-(d+1)$, and every output
factors through the two six-dimensional variables $(D,A)$;
hence the rank is twelve. Multiplication of the character
by $e^{b\theta}$ conjugates this calculation by the invertible
shear $\eta'=\eta+bA+b^2B$, $A'=A+2bB$.
This verifies the convention needed for the rational
polarization on the quotient. The complete convention
verification is retained in the working report.

It remains to verify the actual object. Write
$D_0=\bar\nu_*\mathcal O_{\widetilde Z/G}$.
It is a pure sheaf of dimension two: it injects into the
finite pushforward of the structure sheaves of the smooth
normalization components. The map $\mathcal O_Y\to D_0$
has kernel $I=\mathcal I_{\bar Z_{\mathrm{surf}}}$ and
zero-dimensional cokernel $Q$. Thus $\bar E$ has cohomology
$I$ in degree zero and $Q$ in degree one.
Since $I$ is torsion free,
$\operatorname{Hom}(Q,I)=0$. The truncation triangles show
that this implies
$\operatorname{Ext}^{j}(\bar E,\bar E)=0$ for $j<0$.
They also give
$\operatorname{Hom}(\bar E,D_0[-1])=0$, because
$\operatorname{Hom}(Q,D_0)=0$ by purity.
Codimension two of $\operatorname{Supp}D_0$ on the smooth
fourfold gives
$\operatorname{Hom}(D_0,\mathcal O_Y)
=\operatorname{Ext}^1(D_0,\mathcal O_Y)=0$.
Applying $\operatorname{Hom}(-,\mathcal O_Y)$ to
$\bar E\to\mathcal O_Y\to D_0$ therefore yields
$\operatorname{Hom}(\bar E,\mathcal O_Y)=\mathbf C$.
Applying $\operatorname{Hom}(\bar E,-)$ and using the other
vanishing embeds $\operatorname{End}(\bar E)$ into this
one-dimensional space. The identity is nonzero, so
$\bar E$ is simple.

Normalizing a transverse meeting point of two smooth
surfaces adds one skyscraper class to the structure sheaf.
The chosen four points for each of the $\binom\ell2$ pairs
therefore add $4\binom\ell2=2d\ell$ point classes upstairs.
This is exactly the number of isolated points in S-19.
Consequently the partial-normalization complex has the same
K-character as that first ideal seed, as asserted in
Markman Example~8.2.3. The quotient has
$\chi(\bar E,\bar E)=8d$, virtual rank one, and character
$e^{-\theta}w$ in the rational quotient polarization.
The preceding linear calculation gives
$\operatorname{rank}c_{\bar E}^2=12$.
The unit--trace composite is again rank one, so it injects
$H^1(\mathcal O_Y)$ into
$\operatorname{Ext}^1(\bar E,\bar E)$.
The negative-Ext vanishing, Serre duality, and simplicity
now give
\[
 \operatorname{ext}^2(\bar E,\bar E)
 =8d-2+2\operatorname{ext}^1(\bar E,\bar E)
 \ge8d+6>12.
\]
This contradicts (S20.1) if semiregularity is assumed.

For the quantitative assertion, put
$m=\operatorname{ext}^2(\bar E,\bar E)$ and
$a=\operatorname{rank}\sigma_{\bar E}
=\operatorname{rank}\operatorname{ob}_{\bar E}$.
The rank inequality for the composition gives
$12\ge2a-m$. Hence
\[
 \dim\ker\sigma_{\bar E}=m-a
 \ge\frac{m-12}{2}\ge4d-3.
\]
The argument proving simplicity also applies upstairs.
All linearizations therefore differ by characters of $G$,
and their descents by torsion line bundles, as in S-19.
Tensor equivalence preserves the nonzero full kernel,
proving the weak-equivariant assertion for this action.
\end{proof}

For both S-19 and S-20, taking an external product with a
nonzero perfect object preserves the displayed kernel:
the functor sends $u$ to $u\boxtimes\operatorname{id}$,
which is nonzero by K\"unneth, and the full trace is functorial.
The same is true after a derived equivalence.
Thus the canonical product and Orlov-transform route does not
repair these particular factors. This statement does not
exclude a different object or a different finite target action.

\begin{proposition}[S-21: All finite target actions obey the fourfold secant obstruction]
\label{prop:semireg-all-finite-target-actions}
Let $A$ be a complex abelian fourfold, $\mathcal C=
\operatorname{Perf}(A)$, and let a finite subgroup
$G\subset A\times\widehat A$ have a specified coherent action
on $\mathcal C$ by the corresponding translations and
tensor products with degree-zero line bundles. In particular,
no nonidentity group element represents the pair $(0,\mathcal O_A)$.
For a lift $\widetilde E\in\mathcal C^G$ of a perfect object
$E\in\mathcal C$, suppose
$\operatorname{Ext}^{j}_{\mathcal C^G}(\widetilde E,
\widetilde E)=0$ for $j<0$. Suppose further that
\[
 e^B\operatorname{ch}(E)=a\alpha+b\beta,
 \quad
 \alpha=\frac{e^{s\theta}+e^{-s\theta}}2,
 \quad
 \beta=\frac{e^{s\theta}-e^{-s\theta}}{2s},
 \quad s^2=-d,
 \tag{S21.1}
\]
where $d>0$, $\theta$ is a rational ample class,
$B$ is a rational algebraic divisor class, and
$a,b\in\mathbf Q$ are not both zero.
If $\widetilde E$ is semiregular in $\mathcal C^G$, then
\[
 \operatorname{ext}^2_{\mathcal C^G}(\widetilde E,
              \widetilde E)=12,
 \qquad
 \operatorname{ext}^0_{\mathcal C^G}(\widetilde E,
              \widetilde E)
 \ge 2+\frac{\chi_A(E,E)}{2|G|}>2,
 \tag{S21.2}
\]
where
\[
 \chi_A(E,E)=8d(da^2+b^2)\frac{\int_A\theta^4}{24}>0.
 \tag{S21.3}
\]
Consequently no such lift that is simple in the equivariant
category is semiregular. In particular, an underlying simple
perfect object with (S21.1) and no negative self-Ext admits no
weakly semiregular lift for any such finite target action.
The conclusion allows arbitrary mixtures of translations
and degree-zero tensor operations; it does not require a
Fourier--Mukai model on which the whole group acts freely
by translations.
\end{proposition}
\begin{proof}
Write $\phi_g$ for the functor associated to $g=(x,P)$.
The invariant category is smooth and proper in characteristic
zero; see Perry 2026, \S5, opening paragraph, p.~12,
and Bayer--Perry Proposition~3.15, author PDF pp.~20--21.
First we verify that all nonidentity coefficient sectors vanish:
\[
 HH^*(\mathcal C,\phi_g)=0\qquad(g\ne1).
 \tag{S21.4}
\]
The coefficient space is the derived Hom from the diagonal
kernel to the kernel on the graph of translation by $x$,
with its line-bundle coefficient. This is the kernel
description in Perry, \emph{Hochschild cohomology and group
actions}, arXiv:1807.09268v2, Remark~4.1 and the proof of
Corollary~4.8, pp.~9 and~12. If $x\ne0$, the graph is
disjoint from the diagonal, so the derived Hom is zero.
If $x=0$ and $g\ne1$, then $P\ne\mathcal O_A$ and
the sheafified HKR decomposition identifies the coefficient
groups with sums of
$H^q(A,\wedge^pT_A\otimes P)$.
They vanish, because $T_A$ is trivial and a nontrivial
degree-zero line bundle on an abelian variety has no
cohomology. The latter also follows from the transform of
$\mathcal O_A$ in Mukai Theorem~2.2, p.~156, by derived
base change away from the origin.

Translations and degree-zero tensor operations act trivially
on $HH^*(A)$ and $HH_*(A)$: under HKR the former fix invariant
forms and vector fields, and the latter have zero Atiyah
class and Chern character one. Perry's Theorem~4.4 and
Remark~4.5, pp.~10--12, now give
\[
 HH^*(\mathcal C^G)\cong HH^*(\mathcal C).
 \tag{S21.5}
\]
The identity-sector inclusion is compatible with the
Hochschild algebra and with evaluation on equivariant objects.
This conclusion uses the specified coherent action, not
merely an unlinearized list of isomorphic functors.

The invariant category has Serre functor $[4]$.
To check this directly, the ordinary Serre trace pairing on
$\operatorname{Perf}(A)$ is invariant under translations
and degree-zero tensor products: translations preserve the
chosen invariant volume form, and tensoring preserves
perfect trace. Conjugation by the linearizations also
preserves trace. Taking $G$-invariants in the Serre duality
isomorphism remains a perfect duality, since $|G|$ is
invertible. This is exactly the invariant-duality argument
in the proof of Perry Lemma~6.5, equations (6.1)--(6.3),
p.~17; the proof there treats the general finite-group
Serre descent before its stated cyclic specialization.

Perry Definition~6.1 and Remark~6.2, p.~16, identify
Hochschild homology with Serre-coefficient Hochschild
cohomology. Converting his cohomological grading to the
usual homological convention, a category with Serre functor
$[4]$ has $HH_j\cong HH^{4-j}$ as vector spaces.
Thus (S21.5) gives equal finite dimensions for
$HH_j(\mathcal C^G)$ and $HH_j(\mathcal C)$, in particular
for $j=0,-2$.

Let $F:\mathcal C^G\to\mathcal C$ be forgetting and
$I:\mathcal C\to\mathcal C^G$ induction. These functors
are mutually left and right adjoint by Perry Lemma~3.3,
p.~7, and $FI\cong\bigoplus_{g\in G}\phi_g$.
Functoriality and additivity of Hochschild homology imply
\[
 F_*I_*=|G|\operatorname{id}_{HH_*(\mathcal C)}.
\]
One may read this additivity from Shklyarov, arXiv:0710.1937v3,
Proposition~3.2 and Theorem~3.4, pp.~14--16, by applying
the latter to a sum of functor kernels.
Equal dimensions then show that $F_*$ is an isomorphism
in degrees zero and minus two, and that
\[
 I_*F_*=|G|\operatorname{id}_{HH_j(\mathcal C^G)}
       \qquad(j=0,-2).
 \tag{S21.6}
\]
In particular, forgetting cannot hide a semiregularity
kernel in a nonidentity Hochschild sector.

We also need the exact Euler scaling, not an estimate by
averaging dimensions. Hochschild Chern characters commute
with functors, so (S21.6) implies
$\operatorname{ch}_{HH}(IF\widetilde E)
=|G|\operatorname{ch}_{HH}(\widetilde E)$.
The categorical Riemann--Roch formula makes the Euler
pairing linear in these characters: its precise source is
Shklyarov Theorem~3.5, p.~17, and its Hom formulation
equations (3.7) and (3.9), pp.~18--19.
The formula applies here via a proper DG-algebra model:
induction of a generator of $\operatorname{Perf}(A)$ is
a generator of $\mathcal C^G$, by adjunction and the
averaging retraction; its endomorphism algebra is proper.
Using the right adjunction $F\dashv I$, we obtain
\[
 \chi_A(E,E)
 =\chi_{\mathcal C^G}(\widetilde E,IF\widetilde E)
 =|G|\chi_{\mathcal C^G}(\widetilde E,\widetilde E).
 \tag{S21.7}
\]

Write $V_i=\operatorname{Ext}^i_A(E,E)^G$, which is the
degree-$i$ self-Ext in $\mathcal C^G$.
Perry 2026 Remarks~2.3 and~2.5, p.~5, give the commuting
semiregularity diagram for forgetting. Since $F_*$ in
degree minus two is an isomorphism, semiregularity of the
lift is equivalent to injectivity of
$\sigma_E|_{V_2}$.
The ordinary evaluation maps
$\operatorname{ob}_E^i:HH^i(A)\to\operatorname{Ext}^i(E,E)$
have image in $V_i$, because $G$ acts trivially on
$HH^i(A)$ and evaluation is natural.
The fourfold adjunction of S-20 restricts to the perfect
invariant Serre pairing on $V_2$. Consequently, if the
lift is semiregular, evaluation in degree two is onto
$V_2$ and
\[
 \dim V_2=\operatorname{rank}c_E^2.
 \tag{S21.8}
\]
Similarly $\dim V_1\ge\operatorname{rank}c_E^1$.
For clarity, the latter needs only the degree-one trace
identity. For $\gamma:\mathcal O_\Delta\to
\mathcal O_\Delta[1]$, naturality of the same adjunction
unit as in 032 equation (5.23) gives
\[
 (1_E\otimes L\Delta^*\gamma)AH_E
     =AH_E[1]\chi_E(\gamma).
\]
Take perfect trace with coefficients and apply HKR as in
(5.20)--(5.24). This gives
$\sigma_E^1\operatorname{ob}_E^1=c_E^1$,
so the claimed rank bound follows in the invariant image.

We calculate both contraction ranks explicitly. A
$B$-field factor $e^{-B}$ does not change them:
conjugation by $e^B$ replaces the degree-one contraction
operator $\iota_v$ by $\iota_v+\iota_vB\wedge$ and fixes
wedge operators. Since $B$ has type $(1,1)$, this is an
invertible shear of $HT^1(A)=U\oplus T$ and hence of its
exterior algebra, while $e^B$ preserves each Hodge-difference
space on the target.
For $w=a\alpha+b\beta$, the degree-two calculation in S-20
has first two blocks
\[
 aD+bA,\qquad\theta\wedge(bD-adA).
\]
The determinant is $-(a^2d+b^2)\ne0$, giving
$\operatorname{rank}c_E^2=12$.
In degree one, put $L=v\mathbin{\lrcorner}\theta\in U$
for $(\eta,v)\in U\oplus T$.
The first two blocks are
\[
 a\eta+bL,\qquad\theta\wedge(b\eta-adL).
\]
The same determinant is nonzero, and wedge by $\theta$
is injective on $U$. Thus $\operatorname{rank}c_E^1=8$.

Finally, $B$ cancels from $\operatorname{ch}(E)^\vee
\operatorname{ch}(E)$. If
$c_\pm=(a\pm b/s)/2$, then
$c_+c_-=(a^2+b^2/d)/4$.
The top term of $w^\vee w$ is
$2c_+c_-(2s)^4\theta^4/24$, proving (S21.3).
Put $h=\dim V_0$. The assumed negative-Ext vanishing and
invariant Serre duality imply vanishing above degree four,
$\dim V_4=h$, and $\dim V_3=\dim V_1$.
Equations (S21.7)--(S21.8) therefore imply
\[
 12=\dim V_2
 =\frac{\chi_A(E,E)}{|G|}-2h+2\dim V_1
 \ge\frac{\chi_A(E,E)}{|G|}-2h+16.
\]
Rearranging gives (S21.2), including its contradiction
when $h=1$. The proof applies separately to every coherent
linearization; none was assumed canonical.
\end{proof}

The earlier fourfold criterion and its free-translation
version are recorded in rStation paper~37, version~4,
Theorems~A, B and $B_G$. S-21 adds the coefficient-sector
and induction argument for arbitrary effective finite
subgroups of $A\times\widehat A$, and permits simplicity
only after equivariantization. It excludes this ample
imaginary-quadratic secant character for equivariantly
simple seeds; it does not exclude arbitrary corrected
characters on the final abelian eightfold.

\begin{proposition}[S-22: The balanced real-multiplication character has no rank-one torsion-free realization]
\label{prop:semireg-rm-balanced-obstruction}
Let $A$ be an abelian fourfold with a Rosati-invariant action
of a real quadratic field $F$, of multiplicity two at each
real embedding on $H^0(T_A)$. Let $\theta$ be an $F$-compatible
polarization and write $\theta=\theta_1+\theta_2$ for its two
real eigenspace components. Put
\[
 v_q=\left(1-\frac q2\theta_1^2\right)
       \left(1-\frac q2\theta_2^2\right),
 \qquad q>0,
 \qquad D=\frac{\int_A\theta^4}{24}.
 \tag{S22.1}
\]
There is no nonzero effective surface cycle whose class
belongs to $\mathbf R\theta_1^2+\mathbf R\theta_2^2$.
Consequently there is no rank-one torsion-free coherent
sheaf $E$ with $e^B\operatorname{ch}(E)=v_q$, for any
rational algebraic divisor class $B$.

The contraction ranks of $v_q$ are
\[
 \operatorname{rank}c_{v_q}^1=8,
 \qquad \operatorname{rank}c_{v_q}^2=18,
 \qquad \chi(Nv_q,Nv_q)=4Dq^2N^2.
 \tag{S22.2}
\]
If a perfect object with normalized character $Nv_q$,
$N\in\mathbf Z_{>0}$, has a gluable, semiregular, equivariantly simple
lift for an effective coherent finite target action as in
S-21, then
\[
 \frac{4Dq^2N^2}{|G|}\in\{2,4\},
 \qquad |G|\in\{Dq^2N^2,\,2Dq^2N^2\}.
 \tag{S22.3}
\]
Here only integer group orders are possible. In particular,
when $G=1$, $q=1$, and $\theta$ is integral, a simple pure
coherent sheaf with normalized character $Nv_1$ cannot be
fully semiregular, for any $N>0$.
\end{proposition}
\begin{proof}
Let $V_1\oplus V_2=H^0(T_A)$ be the two complex
$F$-eigenspaces, each of dimension two, and put
$W_i=V_i^*$. The invariant form representing $\theta_i$
is supported on $V_i$ and has rank two. In particular
$\theta_i^3=0$.
For every mixed holomorphic two-form
$\omega\in W_1\otimes W_2\subset H^0(\Omega_A^2)$,
\[
 (a\theta_1^2+b\theta_2^2)\wedge\omega\wedge\overline\omega=0
 \tag{S22.4}
\]
for all real $a,b$: the first term has excessive holomorphic
and antiholomorphic degree in $V_1$, and the second in $V_2$.

Suppose $Z=\sum n_jZ_j$, $n_j>0$, were such an effective
cycle. On a resolution of each component $Z_j$, the integral
of $\omega\wedge\overline\omega$ is nonnegative, and vanishes
only when the pullback of $\omega$ vanishes. Equation
(S22.4) and positivity therefore force every mixed two-form
to vanish on every component. At a smooth point of a
component, its two-dimensional tangent plane $T$ consequently
has no mixed component in its Pl\"ucker bivector.
Write that bivector as $u_1+u_2$, with
$u_i\in\wedge^2V_i$. Its square is $2u_1\wedge u_2$.
Since it is decomposable, this square vanishes, so one
$u_i$ is zero. Thus $T=V_1$ or $T=V_2$.
On each irreducible component the Gauss map is therefore
constant, with one of these two values.

Here is an explicit Albanese argument for the required
consequence of that constancy. Let $f:Y\to Z_j\subset A$
be a smooth projective resolution, translated so that a
chosen point maps to zero. Every invariant one-form
annihilating $V_i$ pulls back to zero on the dense smooth
locus and hence on $Y$. The induced homomorphism
$\operatorname{Alb}(Y)\to A$ therefore has tangent image
contained in $V_i$. Its image is a positive-dimensional
abelian subvariety, because $f$ is nonconstant. But $V_i$
cannot contain the tangent space of such a subvariety.
Indeed, choose $u\in F\setminus\mathbf Q$. Its action on
$H_1(A,\mathbf Q)$ is rational, while on the underlying
real eigenspace of $V_i$ it is multiplication by the
irrational number $\sigma_i(u)$. That eigenspace contains
no nonzero rational vector: otherwise a nonzero rational
coordinate of $z$ in $uz=\sigma_i(u)z$ would make
$\sigma_i(u)$ rational. The nonzero rational homology of
an abelian subvariety gives a contradiction.
This proves the assertion about effective cycles, without
assuming that $A$ itself is simple.

For a rank-one torsion-free $E$, its double dual is a line
bundle $L$, and $L/E$ has support of codimension at least
two. Its codimension-two character is the effective
fundamental cycle of its surface support. The degree-one
part of $e^B\operatorname{ch}(E)=v_q$ gives
$c_1(L)=-B$. The degree-two part thus identifies that
effective cycle with
$q(\theta_1^2+\theta_2^2)/2$.
This class is nonzero, since its pairing with $\theta^2$
is $4Dq>0$, contradicting the first part.

For completeness, the ranks in (S22.2) can be checked on
the complex Hodge decomposition, without supposing that
$A$ is a product. It factors as the exterior tensor product
of the two two-dimensional eigenspaces. For
$v_i=1-q\theta_i^2/2$, the contraction ranks in degrees
zero, one and two are respectively $1,4,1$.
In degree one, wedge operators on $H^1(\mathcal O)$ and
contraction operators on $H^0(T)$ have independent images
of dimension two each. In degree two the only nonzero
output is the one-dimensional $H^2(\mathcal O)$ block;
the middle wedge--contraction operators vanish for degree
reasons. Consequently the product has degree-one rank
$4+4=8$ and degree-two rank $1+4\cdot4+1=18$.
The three degree-two images lie in distinct pairs of Hodge
differences on the two eigenspace factors, so their sum is
direct. The $B$-field shear of S-21 preserves these ranks.
Since $v_q^\vee=v_q$, the top term of $v_q^2$ is
$q^2\theta_1^2\theta_2^2$. Also
$\int_A\theta_1^2\theta_2^2=4D$, proving the Euler formula.

The all-group argument in S-21 applies to every character.
For a semiregular simple lift here it gives $e_2^G=18$,
$e_1^G\ge8$, and
\[
 \frac{\chi_A(E,E)}{|G|}=2-2e_1^G+18.
\]
The left side is strictly positive, so it is either $2$ or
$4$, proving (S22.3). If $G=1$, $q=1$, and $\theta$ is
integral, then $D\in\mathbf Z_{>0}$ and
$4DN^2\le4$, forcing $D=N=1$.
A pure rank-one sheaf on the smooth integral fourfold is
torsion free, so the first part rules out this remaining
case.
\end{proof}

For $q=1$, (S22.1) is precisely
$(\alpha+\widetilde\alpha)/2$ in the real-multiplication
space of Markman, arXiv:2509.23079, \S11.2.1, pp.~40--41:
\[
 \frac{\cos(\theta_1+\theta_2)+
             \cos(\theta_1-\theta_2)}2
       =\cos(\theta_1)\cos(\theta_2).
\]
Its numerical eligibility was already recorded in the
real-multiplication remark of rStation paper~37, version~4.
S-22 adds the rank-one effectivity obstruction. It does
not exclude a higher-rank weakly equivariant object, or a
rank-one perfect complex which is not a pure sheaf.

\begin{proposition}[S-23: A nonsimple semiregular control on an abelian fourfold]
\label{prop:semireg-nonsimple-control}
For every abelian fourfold $A$ and point $p\in A$, the
perfect sheaf
$E=\mathcal O_A\oplus\mathcal O_p$ is fully semiregular,
with self-Ext dimensions
\[
 (e_0,e_1,e_2,e_3,e_4)=(3,8,12,8,3),
 \qquad \chi(E,E)=2.
 \tag{S23.1}
\]
Its endomorphism algebra is triangular and contains a
nonzero nilpotent. Hence full semiregularity does not
imply simplicity or a semisimple endomorphism algebra.
\end{proposition}
\begin{proof}
The self-Ext algebras of $\mathcal O_A$ and $\mathcal O_p$
are respectively $H^*(\mathcal O_A)$ and
$\wedge^*T_pA$, the latter by a regular local Koszul
resolution. The cross groups are
\[
 \operatorname{Ext}^i(\mathcal O_A,\mathcal O_p)=
 \begin{cases}\mathbf C&i=0,\\0&i\ne0,\end{cases}
\]
and, by Serre duality, the opposite cross groups are
one-dimensional in degree four only. This proves (S23.1)
and identifies the off-diagonal endomorphism
$\mathcal O_A\to\mathcal O_p$ as a nonzero square-zero
endomorphism of $E$.

The Atiyah class of $\mathcal O_A$ is zero, so its
semiregularity map identifies its degree-two self-Ext with
$H^2(\mathcal O_A)$. For the point, the local Koszul
calculation identifies evaluation of a global bivector
with its value in $\wedge^2T_pA$. Thus evaluation is onto
its six-dimensional degree-two self-Ext.
Contraction with $\operatorname{ch}(\mathcal O_p)=[p]$
identifies this bivector space with
$H^4(\Omega_A^2)$, also of dimension six; the other
Hochschild degree-two blocks annihilate $[p]$ for degree
reasons. By 032 Proposition~5.1, equation (5.9), the full
semiregularity map of the point is therefore injective,
with image exactly this last Hodge block.
There are no cross degree-two self-Ext groups for $E$.
Perfect trace on its two diagonal summands has the two
images just described, which are disjoint Hodge blocks.
Their direct sum is injective. This proves full
semiregularity of $E$.
\end{proof}

\begin{proposition}[S-24: Iterated self-extensions do not repair the positive-Euler secant seed]
\label{prop:semireg-iterated-extensions}
Let $A$, $G$ and $\mathcal C^G$ be as in S-21, and let
$\mathcal H$ be the heart of a bounded t-structure on
$\mathcal C^G$. Let $\widetilde F\in\mathcal H$ satisfy
$\operatorname{End}_{\mathcal C^G}(\widetilde F)=\mathbf C$,
with underlying normalized imaginary-quadratic secant
character as in (S21.1), and suppose
\[
 k:=\chi_{\mathcal C^G}(\widetilde F,\widetilde F)\ge2.
 \tag{S24.1}
\]
If $\widetilde E\in\mathcal H$ admits a finite filtration
of length $m>0$ whose every successive quotient is
$\widetilde F$, then $\widetilde E$ is not semiregular.
In particular no iterated self-extension of the quotient
ideal in S-19, in its equivariant coherent heart, repairs
its semiregularity failure: there $k=8d$.
\end{proposition}
\begin{proof}
We first prove the precise endomorphism bound. If $X,Y$
have filtrations of lengths $a,b$ with factors
$\widetilde F$, then
\[
 \dim\operatorname{Hom}_{\mathcal H}(X,Y)\le ab.
 \tag{S24.2}
\]
For $X=\widetilde F$, induction on a filtration of $Y$
and left exactness of $\operatorname{Hom}(\widetilde F,-)$
give the bound $b$, since
$\operatorname{End}(\widetilde F)=\mathbf C$.
Then use a short exact sequence
$0\to X'\to X\to\widetilde F\to0$ and left exactness
of $\operatorname{Hom}(-,Y)$ to obtain
$\dim\operatorname{Hom}(X,Y)\le b+(a-1)b$.
This proves (S24.2), hence $e_0^G(\widetilde E)\le m^2$.

In K-theory $[\widetilde E]=m[\widetilde F]$.
Bilinearity of the Euler pairing gives
$\chi_G(\widetilde E,\widetilde E)=m^2k$.
Its character is multiplied by $m$, which does not
change either contraction rank, namely $8$ and $12$.
Objects of a heart have no negative self-Ext.
If $\widetilde E$ were semiregular, S-21 would therefore
give $e_2^G=12$, $e_1^G\ge8$, and invariant Serre duality
would imply
\[
 12=m^2k-2e_0^G+2e_1^G
 \ge m^2(k-2)+16\ge16,
\]
a contradiction. No splitting of the filtration, and no
injection of the original trace kernel into that of an
extension, was assumed.
\end{proof}

The hypothesis (S24.1) is essential to this estimate.
S-24 concerns a fixed factor with scalar endomorphisms and a fixed coherent
action; it does not exclude arbitrary cones, constituents
with different numerical characters, or a newly chosen
finite action not preserving such a filtration.

\begin{proposition}[S-25: Correlated finite target actions do not repair a simple secant product]
\label{prop:semireg-correlated-products}
Let $X$ be an abelian fourfold and $Y$ an abelian variety.
Let $F\in\operatorname{Perf}(X)$ and
$K\in\operatorname{Perf}(Y)$ satisfy
\[
 \operatorname{End}(F)=\operatorname{End}(K)=\mathbf C,
 \qquad \operatorname{Ext}^{j}(F,F)=
        \operatorname{Ext}^{j}(K,K)=0\quad(j<0).
 \tag{S25.1}
\]
Assume that $F$ has the normalized ample
imaginary-quadratic secant character of (S21.1).
For any derived equivalence
$\Phi:\operatorname{Perf}(X\times Y)\simeq
\operatorname{Perf}(B)$ to an abelian variety $B$, put
$E=\Phi(F\boxtimes K)$.
Then $E$ admits no semiregular lift for any effective
finite subgroup $G\subset B\times\widehat B$ with a
specified coherent translation--tensor action.
The group need not be a product of factor actions.

In particular, changing the finite target action cannot
make the actual genus-four secant-square objects
$\Phi(F_1\boxtimes F_2^\vee)$ weakly semiregular when the
factors are simple and gluable. The assertion also applies
when the second factor has a different Chern character;
only the first factor is required to lie in the ample
imaginary-quadratic secant plane.
\end{proposition}
\begin{proof}
We spell out the coherence issue, since projective factor
linearizations need not inherit the given product cocycle.

First, the identity-sector argument of S-21 works on an
abelian variety $Z$ of any dimension $g$. For an effective
finite translation--tensor action $H$, all nonidentity
Hochschild coefficient sectors vanish, and $H$ acts
trivially on the identity sector. The invariant category
is smooth proper with Serre functor $[g]$. Perry's
Hochschild-cohomology formula and the identification
$HH_j\cong HH^{g-j}$ give equality of homology dimensions.
Induction and forgetting satisfy
$\operatorname{Forg}_*\operatorname{Ind}_*=|H|$.
Consequently
\[
 HH_{-2}(\operatorname{Perf}(Z)^H)
   \longrightarrow HH_{-2}(Z)
 \quad\hbox{is an isomorphism}.
 \tag{S25.2}
\]
Thus a nonzero invariant element of the ordinary full
trace kernel also gives a nonzero equivariant trace kernel.
This uses the same primary inputs as S-21 and is
independent of any choice of scalar multiplication cocycle.

Next, a finite subgroup of $Z\times\widehat Z$ stabilizing
a simple gluable object $Q$ can be given a coherent action
for which $Q$ has a lift. Choose representative functors
$\phi_h$ and isomorphisms $\lambda_h:\phi_hQ\to Q$,
with the identity choices normalized. Natural isomorphisms
$\phi_g\phi_h\to\phi_{gh}$ form a scalar torsor,
since $HH^0(Z)=\mathbf C$. Their evaluations at $Q$ also
form a scalar torsor, since $\operatorname{End}(Q)=\mathbf C$,
and evaluation is an isomorphism between those torsors.
There is therefore a unique multiplication isomorphism
$\mu_{g,h}$ satisfying
\[
 \lambda_{gh}(\mu_{g,h})_Q
       =\lambda_g\,\phi_g(\lambda_h).
 \tag{S25.3}
\]
Both associativity composites agree after evaluation at
$Q$, by (S25.3) and naturality. They consequently agree as natural
transformations, since evaluation is faithful on scalars.
This produces a unital two-categorical action and a
compatible lift of $Q$.
There is no missing higher coherence: $HH^j(Z)=0$ for
$j<0$, and $\operatorname{Ext}^j(Q,Q)=0$ for $j<0$.
Bayer--Perry Lemma~3.3 and Corollary~3.4, pp.~15--16,
identify two-categorical lifts with enhanced lifts in
this situation. Equivalently, the spaces of natural
isomorphisms and of isomorphisms of the framed object
have no higher homotopy groups; their scalar torsors
cancel in the framed automorphism groupoid.
Changing any $\lambda_h$ or the multiplication by
scalars does not change its conjugation action on
$\operatorname{Ext}^*(Q,Q)$.

By Orlov Theorem~2.10 and Corollary~2.13, author PDF
pp.~9--12, conjugating by $\Phi$ identifies
$B\times\widehat B$ with
$(X\times Y)\times\widehat{(X\times Y)}$ and sends
translation--tensor operations to translation--tensor
operations. Transport the given action and any proposed
lift of $E$ across $\Phi$. We obtain an effective finite
subgroup
\[
 H\subset (X\times\widehat X)\times(Y\times\widehat Y)
\]
with a lift of $F\boxtimes K$. Let $H_X,H_Y$ be its
projections. Every projected functor stabilizes its
corresponding factor. Here is a proof valid even for
perfect complexes, without a common-heart assumption.
For $h\in H$ write its component functors as $\phi$ and
$\psi$, so
\[
 \phi F\boxtimes\psi K\simeq F\boxtimes K.
 \tag{S25.4}
\]
Choose a point of $Y$ at which the derived fiber of $K$
is nonzero. Restriction of (S25.4) gives an isomorphism
between finite direct sums of shifts of $\phi F$ and
finite direct sums of shifts of $F$, with both sums
nonzero. A diagonal component of the composite of this
isomorphism and its inverse is the identity of a shift
of $F$. Since its endomorphisms are scalar, at least one
summand of this composite is a nonzero scalar. It makes
that shift of $F$ a retract of a shift of $\phi F$.
The latter also has scalar endomorphisms, so the retraction
is an isomorphism. Thus $\phi F\simeq F[k]$ for some
integer $k$. Translations and degree-zero tensor products
preserve the least and greatest nonzero standard
cohomological degrees, so boundedness and nonvanishing
force $k=0$. The same argument with the factors exchanged
shows $\psi K\simeq K$.

Use (S25.3) to give $H_X$ a coherent action and a lift of
$F$. This group is effective, and the lift has scalar
endomorphisms. S-21 therefore implies that it is not
semiregular. By (S25.2), there exists
\[
 0\ne u\in
 \bigl(\ker(\sigma_F:\operatorname{Ext}^2(F,F)
                         \to HH_{-2}(X))\bigr)^{H_X}.
 \tag{S25.5}
\]
The action on this space is the canonical conjugation
action determined by the stabilizing functor classes;
it is unchanged by the scalar choices used to arrange
coherence.

The class $u\boxtimes\operatorname{id}_K$ is nonzero,
by K\"unneth. It is in the full trace kernel of
$F\boxtimes K$, by naturality under the exact enhanced
functor $(-)\boxtimes K$ (Perry 2026 Remark~2.3).
It is also invariant under the originally given $H$-lift.
Indeed, after the factor stabilizations just proved,
K\"unneth and the negative-Ext vanishing in (S25.1) give
\[
 \operatorname{Hom}(\phi F\boxtimes\psi K,F\boxtimes K)
  =\operatorname{Hom}(\phi F,F)\otimes
       \operatorname{Hom}(\psi K,K),
\]
a one-dimensional space. The given product linearization
therefore differs from the tensor product of any two
factor isomorphisms only by a scalar. Conjugation fixes
$u$ in the first factor and the identity in the second;
the scalar cancels. This addresses correlated actions
and mutually cancelling projective factor cocycles.

Finally, apply (S25.2) on $X\times Y$. The nonzero
invariant class is a nonzero kernel in the equivariant
category, so the given lift is not semiregular. Derived
equivalence preserves this conclusion by full trace
naturality. The proof treats every possible coherent
finite target action and lift separately.
\end{proof}

\begin{proposition}[S-26: An exact tangent-dimension and amplitude gate for the balanced RM character]
\label{prop:semireg-rm-tangent-gate}
In the setting of S-22, let $\widetilde E$ be an
equivariantly simple gluable lift with normalized
character $Nv_q$, and set
$k=4Dq^2N^2/|G|$.
Then full semiregularity is equivalent to
\[
 \operatorname{ext}^1_G(\widetilde E,\widetilde E)
        =10-\frac{k}{2}.
 \tag{S26.1}
\]
The positive-Euler possibilities are exactly $k=4$ with
$e_1^G=8$, or $k=2$ with $e_1^G=9$.

There is also a cohomological restriction independent of
semiregularity. A rank-one perfect object with normalized
character $v_q$ cannot have all the following properties:
$H^j(E)=0$ for $j<0$; $H^0(E)$ is rank-one torsion free;
all $H^j(E)$ for $j>0$ have support of codimension at least
two; and those with positive even $j$ have support of
codimension at least three. In particular, the natural
amplitude $[0,1]$ form with torsion-free $H^0$ and
codimension-two-or-more $H^1$ cannot realize this character.
\end{proposition}
\begin{proof}
S-22 gives contraction rank eighteen in degree two.
The restricted CY4 adjunction and the identity
$\sigma\operatorname{ob}=c^2$ imply that the full trace
is injective if and only if $e_2^G=18$: one implication
is adjunction, and for the other the composition already
has rank eighteen on an eighteen-dimensional target
of evaluation. Invariant Serre duality gives
$e_2^G=k-2+2e_1^G$, proving (S26.1).
The rank-eight degree-one image and positivity of $k$
give the two possibilities in S-22.

For the second assertion put $F=H^0(E)$ and
$L=F^{**}$. The higher cohomology does not change rank
or first Chern class, so $c_1(L)=-B$.
The degree-two part of $e^B\operatorname{ch}(F)$ is the
negative of the effective surface cycle of $L/F$.
The only other degree-two contributions are the negatives
of the effective leading cycles of the odd positive-degree
cohomology sheaves; by hypothesis the even positive-degree
terms contribute zero. Thus
$q(\theta_1^2+\theta_2^2)/2$ would be an effective
surface class, contrary to S-22.
\end{proof}

\begin{proposition}[S-27: A two-sheaf cone gives a concrete gluable derived candidate]
\label{prop:semireg-two-sheaf-cone}
Let $A,T$ be nonzero coherent sheaves on a smooth projective
variety, with
\[
 \operatorname{End}(A)=\operatorname{End}(T)=\mathbf C,
 \qquad \operatorname{Ext}^i(T,A)=0\quad(0\le i\le2).
 \tag{S27.1}
\]
For every nonzero $e\in V:=\operatorname{Ext}^3(T,A)$,
form the triangle
\[
 A\longrightarrow E_e\longrightarrow T[-2]
       \xrightarrow{e}A[1].
 \tag{S27.2}
\]
Then $E_e$ is simple and gluable, with
$H^0(E_e)=A$, $H^2(E_e)=T$, and all other cohomology zero.
Its character is $\operatorname{ch}(A)+\operatorname{ch}(T)$,
and
\[
 V/\mathbf C e\hookrightarrow
       \operatorname{Ext}^1(E_e,E_e).
 \tag{S27.3}
\]
Consequently, in the nonequivariant principal balanced-RM
case of S-26, such a cone with the required character is
fully semiregular precisely when
$\operatorname{ext}^1(E_e,E_e)=8$; in particular it requires
$\dim V\le9$.
\end{proposition}
\begin{proof}
Put $B=T[-2]$. There are no negative self-Ext groups of
$A$ or $B$. Also
$\operatorname{Hom}(A,B[j])=\operatorname{Ext}^{j-2}(A,T)=0$
for $j<0$, while
$\operatorname{Hom}(B,A[j])=\operatorname{Ext}^{j+2}(T,A)=0$
for $j<0$ by (S27.1). Taking extensions preserves these
negative-Hom vanishings, by the long exact Hom sequences
in both variables. Thus $E_e$ is gluable. The cohomology
and character assertions follow directly from (S27.2).

Moreover $\operatorname{Hom}(A,B)=
\operatorname{Hom}(A,B[-1])=0$, so
$\operatorname{Hom}(A,E_e)=\operatorname{End}(A)=\mathbf C$.
Similarly $\operatorname{Hom}(E_e,B)=\operatorname{End}(B)
=\mathbf C$. Every endomorphism of $E_e$ therefore induces
scalar endomorphisms $a,b$ on the two ends of the triangle.
The obstruction to this morphism of triangles is
$(a-b)e$, so $a=b$. The kernel of the map to the two
end endomorphisms is a quotient of
$\operatorname{Hom}(B,A)=\operatorname{Ext}^2(T,A)=0$.
Hence $\operatorname{End}(E_e)=\mathbf C$.

For (S27.3), write $i:A\to E_e$ and $p:E_e\to B$
for the triangle maps. Define
\[
 j(v)=i[1]\,v\,p\in\operatorname{Hom}(E_e,E_e[1]),
 \qquad v\in V=\operatorname{Hom}(B,A[1]).
\]
The long exact sequence obtained by applying
$\operatorname{Hom}(E_e,-)$ to the shifted triangle
contains
\[
 \operatorname{Hom}(E_e,B)\xrightarrow{e\circ(-)}
 \operatorname{Hom}(E_e,A[1])\xrightarrow{i[1]\circ(-)}
 \operatorname{Hom}(E_e,E_e[1]).
\]
Its first map is zero, since $\operatorname{Hom}(E_e,B)=
\mathbf C p$ and $ep=0$. Hence the second map is injective.
Applying $\operatorname{Hom}(-,A[1])$ to (S27.2) gives
\[
 \operatorname{End}(A[1])\xrightarrow{(-)\circ e}
 V\xrightarrow{(-)\circ p}\operatorname{Hom}(E_e,A[1]).
\]
The kernel of the latter map is exactly $\mathbf C e$.
Thus $\ker j=\mathbf C e$, proving (S27.3) directly.
The last assertion follows from S-26.
\end{proof}

S-27 is an explicit construction conditional on its two
coherent sheaves and the character equation. No such pair
with the nonsplit RM character and the required tangent
dimension has been constructed here. It isolates an
allowed amplitude $[0,2]$ and a computable extension space,
in contrast to the excluded amplitude $[0,1]$ forms.

\begin{proposition}[S-28: The two-sheaf cone has an exact Yoneda tangent criterion]
\label{prop:semireg-cone-exact-tangent}
Retain the hypotheses and notation of S-27, and put
$B=T[-2]$. Define
\[
 \delta_e:\operatorname{Ext}^1(A,A)\oplus
                 \operatorname{Ext}^1(B,B)
        \longrightarrow\operatorname{Ext}^2(B,A),
 \qquad
 \delta_e(a,b)=a[1]e+e[1]b.
 \tag{S28.1}
\]
Here the sign is the usual convention in which the shifted
triangle has first two maps $i[1],p[1]$ and last map
$-e[1]$. There is a natural exact sequence
\[
 0\longrightarrow V/\mathbf C e
   \xrightarrow{j}\operatorname{Ext}^1(E_e,E_e)
   \xrightarrow{r}\ker\delta_e\longrightarrow0.
 \tag{S28.2}
\]
In particular,
\[
 \operatorname{ext}^1(E_e,E_e)
   =\dim V-1+\dim\ker\delta_e.
 \tag{S28.3}
\]

Suppose now that the ambient variety is an abelian fourfold
in the principal nonequivariant balanced-RM case of S-26,
and that $e^{\mathfrak b}\operatorname{ch}(E_e)=v_1$,
where $\mathfrak b$ denotes the normalizing $B$-field.
Let
\[
 \rho:HH^1(X)\longrightarrow
       \operatorname{Ext}^1(A,A)\oplus
       \operatorname{Ext}^1(T,T)
 \tag{S28.4}
\]
be simultaneous Hochschild evaluation, equivalently the
infinitesimal translation--tensor action on the two sheaves,
and put $h=\dim\ker\rho$. Then
\[
 E_e\text{ is fully semiregular}
 \quad\Longleftrightarrow\quad
 \dim V=h+1\ \text{ and }\ \ker\delta_e=\operatorname{im}\rho.
 \tag{S28.5}
\]
If the simultaneous infinitesimal stabilizer is zero,
the remaining conditions are exactly
\[
 \dim\operatorname{Ext}^3(T,A)=1,
 \qquad
 \operatorname{rank}\delta_e=
       \operatorname{ext}^1(A,A)+\operatorname{ext}^1(T,T)-8.
 \tag{S28.6}
\]
No existence of sheaves satisfying these conditions and
the balanced-RM character equation is asserted.
\end{proposition}
\begin{proof}
For coherent $A,T$ one has
$\operatorname{Hom}(A,B[k])=0$ for $k\le1$.
The long exact Hom sequences therefore give isomorphisms
\[
 \operatorname{Ext}^1(A,A)
     \xrightarrow{i[1]\circ(-)}\operatorname{Hom}(A,E_e[1]),
 \qquad
 \operatorname{Ext}^1(B,B)
     \xrightarrow{(-)\circ p}\operatorname{Hom}(E_e,B[1]).
 \tag{S28.7}
\]
Consequently every $u:E_e\to E_e[1]$ determines unique
$a,b$ by $ui=i[1]a$ and $p[1]u=bp$; set $r(u)=(a,b)$.
The morphism-of-triangles axiom and uniqueness of $b$
imply $a[1]e=-e[1]b$. Conversely, if this equality holds,
apply the same axiom to the rotated triangles, starting
with the commuting square on $B\to A[1]$, to obtain a
map $u$ inducing $a,b$. Thus $\operatorname{im}r=\ker\delta_e$.

If $r(u)=0$, then $ui=0$, so exactness gives $u=wp$ for
some $w:B\to E_e[1]$. Since $p[1]u=0$ and precomposition
with $p$ is injective in the second map of (S28.7),
$p[1]w=0$. A further long exact Hom sequence gives
$w=i[1]v$ for some $v\in V$. Hence $u=j(v)$.
The converse is immediate from consecutive triangle maps
having zero composite. S-27 already proves
$\ker j=\mathbf C e$, so (S28.2) and (S28.3) follow.

For the abelian-fourfold assertion, write
$\operatorname{ob}^1_E:HH^1(X)\to\operatorname{Ext}^1(E_e,E_e)$.
The degree-one contraction rank is eight by S-22.
Since $\sigma^1\operatorname{ob}^1_E=c_E^1$ and
$\dim HH^1(X)=8$, the map $\operatorname{ob}^1_E$ is
injective. Naturality gives
$r\operatorname{ob}^1_E=\rho$; the even shift $[-2]$
does not change the identification for the $T$ term.
Thus $\operatorname{im}\rho\subseteq\ker\delta_e$,
$\dim\operatorname{im}\rho=8-h$, and
\[
 \dim\bigl(\operatorname{im}\operatorname{ob}^1_E
                 \cap\ker r\bigr)=h\le\dim V-1.
 \tag{S28.8}
\]
If $E_e$ is fully semiregular, S-26 gives
$\operatorname{ext}^1(E_e,E_e)=8$. From (S28.3) and
$\dim\ker\delta_e\ge8-h$ one gets $\dim V-1\le h$.
Together with (S28.8), this forces $\dim V=h+1$ and
$\ker\delta_e=\operatorname{im}\rho$.
Conversely those two equalities make (S28.3) equal to
eight, so S-26 gives full semiregularity. Setting $h=0$
proves (S28.6).
\end{proof}

\begin{proposition}[S-29: The transverse-surface cone requires special cohomology]
\label{prop:semireg-surface-cone-cohomology}
Let $X$ be a simple complex abelian fourfold, and let
$Z,S\subset X$ be smooth connected surfaces meeting
transversely in the finite scheme $D$. Write $i:S\hookrightarrow X$,
let $L$ be a line bundle on $X$ and $M$ a line bundle on $S$,
and set
\[
 A=I_Z\otimes L,\qquad T=i_*M,\qquad
 M'=M\otimes L^{-1}|_S,\qquad
 N=K_S\otimes(M')^{-1},\qquad d=\operatorname{length}D.
\]
There are natural identifications
\[
 \operatorname{Ext}^{j+2}_X(T,A)=H^j(S,N\otimes I_{D/S})
       \quad(0\le j\le2),
 \qquad
 \operatorname{Ext}^{0,1}_X(T,A)=0.
 \tag{S29.1}
\]
The two sheaves have scalar endomorphisms, and each
translation--tensor evaluation $HH^1(X)\to\operatorname{Ext}^1$
is injective.

Suppose that a cone of S-27 formed from these sheaves has
the principal balanced-RM normalized character $v_1$.
If it is fully semiregular, then
\[
 \begin{split}
 H^0(S,N\otimes I_D)&=0,\\
 h^1(S,N\otimes I_D)&=1,\\
 h^0(S,M')&\ge8,\\
 h^2(S,M')-h^1(S,M')&=d-1\ge1.
 \end{split}
 \tag{S29.2}
\]
In fact the only possibilities are
\[
 (h^1(S,M'),h^2(S,M'))=(0,d-1)\quad\text{or}\quad(1,d).
 \tag{S29.2a}
\]
In the first case evaluation of $N$ at $D$ has
one-dimensional cokernel; in the second it is an
isomorphism. Together with the full character equation, the first two
conditions and the rank condition (S28.6) are sufficient
as well. In particular, choosing $M'$ with
$H^{>0}(S,M')=0$ cannot give a fully semiregular cone.
More precisely, for every such acyclic choice satisfying
the character equation and every nonzero extension,
\[
 \dim\ker\bigl(\sigma^2_{E_e}\bigr)\ge d-1\ge1.
 \tag{S29.3}
\]
\end{proposition}
\begin{proof}
Transversality gives
$Li^*I_Z=I_{D/S}$: pull back the ideal triangle and
compute the other term by the Koszul complex of $Z$,
which remains regular on $S$.
For the codimension-two regular immersion $i$, duality
gives $i^!\mathcal O_X=\det N_{S/X}[-2]=K_S[-2]$ and
$i^!A=Li^*A\otimes K_S[-2]$.
The precise general statements are Stacks Project,
Lemma~48.9.7 (tag \texttt{0A9X}),
Lemma~48.10.3 (tag \texttt{0E2N}), and
Lemma~48.15.5 (tag \texttt{0BQZ}); the last is the
Koszul duality computation with shift $[-2]$.
Adjunction therefore yields
\[
 R\operatorname{Hom}_X(i_*M,I_Z\otimes L)
     =R\Gamma(S,N\otimes I_D)[-2],
 \tag{S29.4}
\]
which proves (S29.1).
Endomorphisms of $A$ extend to its line-bundle double
dual and hence are scalar; those of $T$ are
$H^0(S,\mathcal O_S)=\mathbf C$.

A translation--tensor operation fixing $A$ must preserve
$Z$, and its degree-zero tensor factor is then determined
by the double dual $L$. The translation stabilizer of
$Z$ is finite: its connected component is an abelian
subvariety of the simple variety $X$, and cannot be $X$.
For $T$, preservation of the support likewise forces the
translation into the finite stabilizer of $S$. With
translation fixed, the remaining stabilizer lies in
$\ker(\operatorname{Pic}^0(X)\to\operatorname{Pic}^0(S))$,
which has zero Lie algebra: $S$ generates $X$ by simplicity,
so its Albanese map onto $X$ is surjective and the dual
map on Picard varieties has injective differential.
Thus both infinitesimal stabilizers vanish, proving the
two rank-eight evaluation assertions.

In particular $h=0$ in S-28. Its criterion gives the
first two conditions of (S29.2) and
\[
 \operatorname{rank}\delta_e
   =\operatorname{ext}^1(A,A)+\operatorname{ext}^1(T,T)-8
   \ge8.
\]
The target of $\delta_e$ has dimension
$h^2(N\otimes I_D)=h^2(N)=h^0(M')$ by Serre duality,
so $h^0(M')\ge8$.
In the cohomology sequence of
$0\to N\otimes I_D\to N\to N|_D\to0$, the first two
conditions imply
\[
 1=d-h^0(N)+h^1(N)
   =d-h^2(M')+h^1(M').
\]
Equivalently, $h^2(M')-h^1(M')=d-1$.
More precisely that exact sequence is
\[
 0\longrightarrow H^0(N)\longrightarrow H^0(N|_D)
 \longrightarrow \mathbf C\longrightarrow H^1(N)
 \longrightarrow0.
 \tag{S29.5}
\]
It follows that $h^1(N)=h^1(M')$ is either zero or one;
the two cases and the assertions about evaluation follow.

For the lower bound on $d$, the first Chern equation
forces the normalizing $B$-field to be $-c_1(L)$.
The codimension-two character equation therefore gives
\[
 [Z]-[S]=C,\qquad C=\frac{\theta_1^2+\theta_2^2}{2}.
 \tag{S29.6}
\]
Both $\theta_i$ have invariant semipositive representatives,
so $C\cdot[S]\ge0$. On a simple complex abelian variety
every irreducible subvariety is geometrically nondegenerate;
Debarre, \emph{Fulton--Hansen and Barth--Lefschetz theorems
for subvarieties of abelian varieties} (1995),
Examples~(1), pp.~188--189, and Proposition~1.4,
pp.~189--190, imply $S-S=X$.
Thus the difference map $S\times S\to X$ is generically
finite and has positive degree $[S]^2$.
This number is even: in
$H^*(X,\mathbf Z)=\bigwedge^*H^1(X,\mathbf Z)$ the square
of every degree-four integral class has even top
coefficient, because each basis monomial squares to zero
and each pair of distinct monomials occurs twice.
Consequently
\[
 d=[Z]\cdot[S]=[S]^2+C\cdot[S]\ge2.
 \tag{S29.7}
\]

Finally, if $H^{>0}(M')=0$, Serre duality gives
$H^0(N)=H^1(N)=0$. The same restriction sequence gives
$H^0(N\otimes I_D)=0$ and
$V=H^1(N\otimes I_D)\simeq H^0(N|_D)$, of dimension $d$.
The first map in the rank-eight simultaneous orbit is
injective, so $\dim\ker\delta_e\ge8$.
S-28 implies $\operatorname{ext}^1(E_e,E_e)\ge d+7$,
and the CY4 Euler equation, with $\chi(E_e,E_e)=4$,
gives $m:=\operatorname{ext}^2(E_e,E_e)\ge2d+16$.
Put $k=\dim\ker\sigma^2_{E_e}$.
By the CY4 transpose identity,
$\operatorname{rank}\sigma^2_{E_e}
 =\operatorname{rank}\operatorname{ob}^2_{E_e}=m-k$,
while the composition has rank eighteen by S-22.
The elementary rank inequality yields
$18\ge2(m-k)-m=m-2k$, hence $k\ge d-1$.
The sufficiency statement is exactly S-28 after (S29.1)
has checked the S-27 hypotheses.
\end{proof}

\begin{theorem}[S-30: An integral lattice obstruction closes the simple rank-one balanced RM gate]
\label{thm:semireg-rm-integral-obstruction}
Let $(X,\theta)$ be a complex abelian fourfold with a
Rosati-invariant action of the real quadratic field
$F=\mathbf Q(\sqrt{\Delta_0})$, where $\Delta_0>1$ is
squarefree, and with the two complex eigenspaces both
of dimension two. Use the notation
\[
 D=\int_X\frac{\theta^4}{24},\qquad
 C=\frac{\theta_1^2+\theta_2^2}{2},\qquad
 v_q=1-qC+\frac{q^2}{4}\theta_1^2\theta_2^2,
 \qquad q\in\mathbf Q_{>0}.
\]
Here $\theta$ is an integral polarization and
$\theta=\theta_1+\theta_2$ is its orthogonal decomposition
along the two eigenspaces.
If $a\in\mathbf Q$ is nonzero and $aC$ is an integral
cohomology class, then
\[
 t=a^2D\in\mathbf Z_{>0},\qquad \Delta_0\mid t.
 \tag{S30.1}
\]
Consequently no rank-one perfect object with normalized
character $v_q$ can be both simple, gluable, and fully
semiregular. More quantitatively, every simple gluable
rank-one perfect object with that normalized character
satisfies
\[
 \dim\ker\sigma_E^2\ge2Dq^2-2\ge2\Delta_0-2>0.
 \tag{S30.2}
\]
In particular, in the principal $q=1$ case, the character
$v_1$ is not the normalized character of any rank-one
perfect object at all. This last assertion does not
require simplicity, gluability, or semiregularity.

For a positive-rank perfect object with
$e^B\operatorname{ch}(E)=Nv_q$, one always has
\[
 N^4Dq^2\in\mathbf Z_{>0},\qquad
 \Delta_0\mid N^4Dq^2.
 \tag{S30.3}
\]
If, in addition, $E$ has a simple gluable fully semiregular
lift for an effective finite coherent target action $G$
as in S-22, write $k=\chi_G\in\{2,4\}$. Then
\[
 \begin{cases}
 \Delta_0\mid N^2|G|,&k=4,\\
 2\Delta_0\mid N^2|G|,&k=2.
 \end{cases}
 \tag{S30.4}
\]
For $G=1$, in particular, $N$ must be a multiple of
$\Delta_0$, and $N$ must also be even in the $k=2$ case.
Moreover, in every rank, a nonequivariant simple gluable
fully semiregular realization requires a genuinely
fractional normalizing class: $B$ cannot be the Chern
class of a line bundle. Equivalently, its determinant
cannot admit an $N$th root in $\operatorname{Pic}(X)$.
The integer $\Delta_0$ here is the squarefree radicand,
not necessarily the discriminant of the quadratic field.
\end{theorem}
\begin{proof}
Put $L=H_1(X,\mathbf Z)$ and let $\nu$ be the primitive
positive generator of $\bigwedge^8L^*$, so
$\theta^4/24=D\nu$. The two real eigenspaces give
$L_{\mathbf R}=U\oplus W$, each of real dimension four.
Write $\nu_U=\theta_1^2/2$ and $\nu_W=\theta_2^2/2$.
Then $C=\nu_U+\nu_W$ and $\nu_U\wedge\nu_W=D\nu$.

Suppose $c=aC\in\bigwedge^4L^*$. Define the integral
endomorphism $R_c$ of $\bigwedge^2L$ by
\[
 \iota_{R_c(x)}\nu=(\iota_xc)\wedge c.
 \tag{S30.5}
\]
Contraction with the primitive volume form is a
unimodular identification between
$\bigwedge^2L$ and $\bigwedge^6L^*$, so $R_c$ is indeed
integral. A direct calculation on the three summands
\[
 \bigwedge^2L_{\mathbf R}
  =\bigwedge^2U\oplus(U\otimes W)\oplus\bigwedge^2W
\]
shows that $R_c=tP$, where $t=a^2D$ and $P$ is the
projector which is the identity on the first and third
summands and zero on the middle summand.
Moreover $c^2=2t\nu$. The square of every integral
four-form on an eight-dimensional lattice has even top
coefficient: basis monomials square to zero and their
cross terms occur in pairs. Hence $t$ is an integer.

Choose the rational endomorphism $T$ of $L_{\mathbf Q}$
given by $\sqrt{\Delta_0}\in F$. It satisfies
$T^2=\Delta_0\operatorname{id}$ and has eigenvalues
$\sqrt{\Delta_0},-\sqrt{\Delta_0}$ on $U,W$. Therefore
\[
 2P-\operatorname{id}
      =\frac{\bigwedge^2T}{\Delta_0}.
 \tag{S30.6}
\]
Let $p$ divide $\Delta_0$. If $p\nmid t$, the integrality
of $tP$ makes $P$ integral over $\mathbf Z_p$.
Thus the left side of (S30.6) is an integral involution,
and hence an automorphism of $\bigwedge^2L_{\mathbf Z_p}$.
Apply Smith normal form over $\mathbf Z_p$ to the
invertible rational matrix $T$. If its elementary-divisor
exponents are integers $m_1,\ldots,m_8$, those of its
exterior square, before reordering, are
$m_i+m_j$ for $i<j$. Equation (S30.6) and unimodularity
would force
\[
 m_i+m_j=v_p(\Delta_0)=1\qquad(i<j).
 \tag{S30.7}
\]
Using any three indices forces every $m_i=1/2$, a
contradiction. This argument works also at $p=2$.
Every prime factor of the squarefree integer $\Delta_0$
therefore divides $t$, proving (S30.1).

For rank one, the degree-two character equation gives
$B=-c_1(E)$. The determinant of a perfect complex is a
line bundle, so $E'=E\otimes\det(E)^{-1}$ is an actual
perfect complex with character $v_q$. Its second Chern
class is $qC$, which is integral. Riemann--Roch on the
abelian fourfold gives $\chi(E')=Dq^2$ and
$\chi(E,E)=4Dq^2$. By (S30.1), $Dq^2$ is a positive
multiple of $\Delta_0$.

If $E$ is simple and gluable, the rank-eight degree-one
contraction forces $e_1\ge8$, and CY4 Serre duality gives
\[
 m:=e_2=4Dq^2-2+2e_1\ge4Dq^2+14.
\]
Put $b=\dim\ker\sigma_E^2$. The full transpose identity
makes the ranks of $\sigma_E^2$ and
$\operatorname{ob}_E^2$ both $m-b$. Their composition has
rank eighteen. Thus
$18\ge2(m-b)-m=m-2b$, which proves (S30.2).
For $D=q=1$, (S30.1) already contradicts $qC$ being
integral; this proves the assertion without simplicity.

For general positive rank $N$, put $\ell=c_1(E)$.
The degree-two equation gives $B=-\ell/N$, and the
next character equation is
\[
 N^2qC=Nc_2(E)-(N-1)\frac{\ell^2}{2}.
 \tag{S30.8}
\]
The class $\ell^2/2$ is integral in the exterior algebra
of a torus. Thus the right side is integral and (S30.1)
with $a=N^2q$ proves (S30.3). In the equivariant case
S-22 gives $k=4Dq^2N^2/|G|\in\{2,4\}$, so
\[
 N^4Dq^2=\frac{kN^2|G|}{4}.
\]
Substitution into (S30.3) gives (S30.4).
For $G=1$, squarefreeness of $\Delta_0$ implies
$\Delta_0\mid N$; when $k=2$, the additional divisibility
by two forces $N$ even.

Finally suppose $G=1$ and $B=c_1(Q)$ for a line bundle
$Q$. Then $E\otimes Q$ is an actual perfect object with
character $Nv_q$, so $s:=\chi(E\otimes Q)=NDq^2$ is a
positive integer. Full semiregularity, simplicity and
gluability give
$4Ns=\chi(E,E)=2-2e_1+18\le4$.
Thus $N=s=1$, which is the excluded rank-one case.
Since $NB=-c_1(E)$, a line-bundle determinant root would
make $B$ integral, and conversely an integral $B$ is a
line-bundle class by the integral $(1,1)$ theorem.
The residual degree-zero line bundle can be absorbed
into an $N$th root because multiplication by $N$ on
$\operatorname{Pic}^0(X)$ is surjective.
\end{proof}

\begin{proposition}[S-31: Rank-two and rank-three bundles cannot realize the balanced character]
\label{prop:semireg-rm-low-rank-bundles}
In the real-multiplication setting of S-30, no vector
bundle of rank $N\in\{1,2,3\}$ has normalized character
$e^B\operatorname{ch}(E)=Nv_q$ with $q>0$.
Consequently a nonequivariant simple gluable fully
semiregular vector-bundle realization would have rank
at least four and divisible by $\Delta_0$.
For example, the smallest remaining ranks allowed by
these two tests are four for $F=\mathbf Q(\sqrt2)$,
and six for $F=\mathbf Q(\sqrt3)$.
These are necessary conditions, not constructions.
\end{proposition}
\begin{proof}
Under the splitting principle, replace the $N$ Chern
roots $x_i$ by $x_i+B$. The resulting normalized Chern
polynomial still has degree at most $N$, regardless of
whether $B$ is the Chern class of a line bundle.
The normalized character has first and third components
zero, second component $-NqC$, and fourth component
$NDq^2\nu$, with $\nu$ the primitive top class.
Newton's identities therefore give
\[
 \bar c_1=\bar c_3=0,\qquad
 \bar c_2=NqC,\qquad
 \bar c_4=\frac{\bar c_2^2-12\overline{\operatorname{ch}}_4}{2}
          =N(N-6)Dq^2\nu.
 \tag{S31.1}
\]
For $N=1$, the nonzero second Chern class contradicts
the degree-one Chern polynomial. For $N=2,3$, the
nonzero fourth class contradicts degree at most $N$.
The final rank restrictions follow from S-30.
\end{proof}

\begin{theorem}[S-32: Two-adic integrality forces an even polarization or group degree]
\label{thm:semireg-rm-two-adic-obstruction}
Let $(X,\theta)$ be a polarized abelian
fourfold with Rosati-invariant real multiplication by
$F=\mathbf Q(\sqrt{\Delta_0})$ of type $(2,2)$, where
$\Delta_0>1$ is squarefree and
$\Delta_0\equiv2$ or $3\pmod4$.
Assume that $H_1(X,\mathbf Z_2)$ is stable under the
maximal local order
$R=\mathbf Z_2[\sqrt{\Delta_0}]$, and that every rational
divisor class has the form
$\ell(x,y)=\theta(\alpha x,y)$ for some $\alpha\in F$.
Put $D=\int_X\theta^4/24$.
The latter hypothesis is the usual minimal divisor space
$\operatorname{NS}(X)_{\mathbf Q}=F\theta$.
If $D$ is odd, then for every integral divisor class
$\ell$, every positive integer $N$, and every rational
two-adic unit $u$, the class
\[
 uC-\frac{\ell^2}{2N},\qquad
 C=\frac{\theta_1^2+\theta_2^2}{2},
 \tag{S32.1}
\]
is not two-adically integral. Consequently, if $D$ is odd,
no positive-rank perfect object has normalized character
$Nv_q$ with $Nq$ a two-adic unit.
More generally, every simple gluable fully semiregular
lift of a character $Nv_q$, $q\in\mathbf Q_{>0}$, for an effective finite coherent
target action $G$ must satisfy
\[
 2\mid D|G|.
 \tag{S32.1a}
\]
In particular, for a principal polarization there is no
nonequivariant simple gluable fully semiregular balanced-RM
perfect object, in any positive rank or for any
normalizing $B$-field.
\end{theorem}
\begin{proof}
For the integral-class assertion assume $D$ is odd.
Write $\omega=\sqrt{\Delta_0}$ and $K=\mathbf Q_2(\omega)$.
The ring $R=\mathbf Z_2[\omega]$ is a discrete valuation
ring: $x^2-\Delta_0$ is Eisenstein when $\Delta_0$ is even;
when $\Delta_0\equiv3\pmod4$, the element $\omega-1$
satisfies the Eisenstein polynomial
$x^2+2x+1-\Delta_0$.
Here we use the Eisenstein DVR criterion, for example
Sutherland, \emph{18.785 Number Theory I}, Fall 2019,
Lecture~11, Lemma~11.4 and its proof, pp.~1--2.
Put $L_2=H_1(X,\mathbf Z_2)$.
It is a finite torsion-free $R$-module, hence free of
rank four (Stacks Project, Lemma~15.22.11,
tag \texttt{0AUW}).
Set $\kappa=1/(2\omega)$.
The trace pairing $(a,b)\mapsto\operatorname{Tr}_{K/\mathbf Q_2}
(\kappa ab)$ on $R$ is perfect: its matrix in the basis
$1,\omega$ is
$\left(\begin{smallmatrix}0&1\\1&0\end{smallmatrix}\right)$.

The self-adjoint $R$-action therefore gives a unique
$R$-bilinear alternating form $\psi:L_2\times L_2\to R$
with
\[
 \theta(x,y)=\operatorname{Tr}_{K/\mathbf Q_2}
                         (\kappa\psi(x,y)).
 \tag{S32.2}
\]
For completeness, define $\psi(x,y)$ by requiring
$\theta(ax,y)=\operatorname{Tr}(\kappa a\psi(x,y))$
for all $a\in R$. Perfectness of the trace pairing gives
existence and uniqueness; self-adjointness gives
bilinearity. Also $\theta(ax,x)=0$ for every $a$, since
it equals both $\theta(x,ax)$ and $-\theta(x,ax)$;
thus $\psi(x,x)=0$ even at two.
The $\psi$-dual lattice equals the $\theta$-dual lattice:
test the trace against every $a\in R$, using $RL_2=L_2$.
Since $D$ is odd, the polarization form is unimodular
over $\mathbf Z_2$ and makes the latter lattice
$L_2$ itself. Hence $\psi$ is perfect.
Elementary symplectic elimination over the local ring
now gives an $R$-basis $e_1,f_1,e_2,f_2$ with
$\psi(e_i,f_j)=\delta_{ij}$ and the other pairings zero.
Indeed choose a primitive vector, pair it to a unit by
perfectness, rescale the second vector, and split off
its orthogonal hyperbolic plane; repeat once.

By the divisor-space hypothesis, $\ell=\theta(\alpha\cdot,\cdot)$
for $\alpha\in F$. Integrality of $\ell$ and unimodularity
of $\theta$ imply $\alpha L_2\subset L_2$.
The free $R$-basis therefore implies
$\alpha=a+b\omega$ with $a,b\in\mathbf Z_2$.
On either hyperbolic plane, the three pairings of $\ell$
corresponding to products $1,\omega,\omega^2$ are
\[
 \operatorname{Tr}(\kappa\alpha)=b,\qquad
 \operatorname{Tr}(\kappa\alpha\omega)=a,\qquad
 \operatorname{Tr}(\kappa\alpha\omega^2)=\Delta_0b.
 \tag{S32.3}
\]
Over a splitting field, $C$ is the sum of the squares
of the two embedding components of $\theta$, divided
by two. Evaluating on one ordered pair from each
hyperbolic plane consequently gives
$\operatorname{Tr}(\kappa^2\omega^j)$ when the product
of the four $R$-coefficients is $\omega^j$.
In particular, on the respective ordered quadruples
\[
 (e_1,f_1,e_2,f_2),\quad
 (\omega e_1,f_1,\omega e_2,f_2),\quad
 (\omega e_1,f_1,e_2,f_2),
\]
the values of $C$ are $1/(2\Delta_0),1/2,0$,
and those of $\ell^2/2$ are $b^2,a^2,ab$.
Thus integrality of (S32.1) would require
\[
 \frac{b^2}{N}-\frac{u}{2\Delta_0}\in\mathbf Z_2,
 \qquad
 \frac{a^2}{N}-\frac{u}{2}\in\mathbf Z_2,
 \qquad
 \frac{ab}{N}\in\mathbf Z_2.
 \tag{S32.4}
\]

Because $u$ is a two-adic unit, if $\Delta_0$ is even,
the first two congruences force
$2v_2(b)-v_2(N)=-2$ and
$2v_2(a)-v_2(N)=-1$, which give incompatible parities
for $v_2(N)$. If $\Delta_0$ is odd, both force valuation
$-1$, so $v_2(a)=v_2(b)$ and
$v_2(ab/N)=-1$, contradicting the third congruence.
This proves the integral-class assertion.

For a perfect object with normalized character
$Nv_q$, put $u=Nq$. The first two Chern equations give
$B=-\ell/N$ and
\[
 c_2(E)=\frac{N-1}{2N}\ell^2+uC
       =\frac{\ell^2}{2}
          +\left(uC-\frac{\ell^2}{2N}\right).
\]
Both $c_2(E)$ and $\ell^2/2$ are integral, contradicting
(S32.1) whenever $u$ is a two-adic unit.
Finally suppose a simple gluable fully semiregular lift
exists and $D|G|$ is odd. By S-22 its invariant Euler
pairing is
$k=4D(Nq)^2/|G|\in\{2,4\}$.
The value $k=2$ would give
$v_2((Nq)^2)=-1$, impossible for a rational square.
The value $k=4$ gives $v_2(Nq)=0$, which has just been
excluded. This proves (S32.1a) and the principal
nonequivariant corollary.
\end{proof}

\begin{theorem}[S-33: Odd ramification obstructs the balanced character in every rank]
\label{thm:semireg-odd-ramification}
Let \((X,\theta)\) be a polarized complex abelian fourfold, with polarization volume \[
 D=\int_X\frac{\theta^4}{24}\in\mathbf Z_{>0}.
\] Let \(F=\mathbf Q(\sqrt d)\) be a real quadratic field, where \(d>1\) is squarefree. Assume its rational action is fixed by the Rosati involution and its two complex eigenspaces each have dimension two. Write \[
 \theta=\theta_++\theta_-,\qquad
 C=\frac{\theta_+^2+\theta_-^2}{2},\qquad
 \nu=\frac{\theta^4}{24}
      =\frac{\theta_+^2\theta_-^2}{4},\qquad \int_X\nu=D.
 \tag{S33.1}
\] For \(r\in F\), let \(\theta_r\) denote the rational alternating form \((x,y)\mapsto\theta(rx,y)\), equivalently \(\sigma_+(r)\theta_++\sigma_-(r)\theta_-\). Fix an odd prime \(p\) ramified in \(F\), so \(p\mid d\), and assume \[
 p\nmid D,\qquad
 (\mathcal O_F\otimes\mathbf Z_p)H_1(X,\mathbf Z_p)
                  \subseteq H_1(X,\mathbf Z_p).
 \tag{S33.2}
\] Thus only maximal-order action at this prime is required.

There is \textbf{no} perfect object \(\mathcal E\) of positive rank \(N\), with \(\ell=c_1(\mathcal E)\in F\theta\), and no \(u\in\mathbf Q^\times\) with \(v_p(u)=0\), satisfying \[
 e^{-\ell/N}\operatorname{ch}(\mathcal E)
       =N-uC+\frac{u^2}{N}\nu.
 \tag{S33.3}
\] The assertion allows an arbitrary fractional normalizing class \(-\ell/N\) and does not require \(\mathcal E\) to be a sheaf, simple or gluable. For the balanced character \(Nv_q\), the scalar here is \(u=Nq\).

In particular, suppose \(\operatorname{NS}(X)_{\mathbf Q}=F\theta\) and maximal-order action is integral at every ramified prime. For a simple gluable fully semiregular lift with an effective finite coherent target action \(G\) as in S-22, every odd ramified prime must divide \(D|G|\). Combining with S-32's even-ramification obstruction gives the necessary condition \[
 \operatorname{rad}(|\operatorname{Disc}(F)|)\mid D|G|.
 \tag{S33.4}
\] In the nonequivariant case, all ramified primes must divide \(D\). In particular, there is no such positive-rank realization for a principal polarization and a genuine real quadratic field. This excludes the proposed normalized character under the stated local maximal-order and determinant hypotheses; it does not assert nonexistence for arbitrary derived characters or arbitrary polarizations/actions.
\end{theorem}
\begin{proof}
\par\medskip\noindent\textbf{1. The integral endomorphism corresponding to the determinant}\quad 

Since \(\ell\in F\theta\), write \(\ell=\theta_r\) for \(r\in F\). Because \(p\nmid D\), the integral alternating form \(\theta\) is perfect over \(\mathbf Z_p\): its determinant is \(D^2\) up to a unit in any integral basis. Integrality of \(\ell\) therefore implies \(rH_1(X,\mathbf Z_p)\subseteq H_1(X,\mathbf Z_p)\). Put \(K=F\otimes\mathbf Q_p\), \(O=\mathcal O_F\otimes\mathbf Z_p\) and \(L=H_1(X,\mathbf Z_p)\). The ring \(O\) is the valuation ring of a ramified quadratic extension. The \(O\)-module \(L\) is torsion-free of rank four, hence free (Stacks Project, Lemma~15.22.11, tag \texttt{0AUW}). The scalar \(r\) preserves \(O^4\), and consequently \(r\in O\).

Set \(s=\sqrt d\). The Eisenstein DVR criterion (Sutherland, Fall 2019 Lecture~11, Lemma~11.4, pp.~1--2) applies to \(x^2-d\). Because \(p\) is odd, \(O=\mathbf Z_p[s]\); the different is the principal ideal \((s)\), and \(s\) is a uniformizer up to a unit. Normalize its valuation by \(v_\pi(s)=1\), so \(v_\pi(p)=2\). We can write \[
 r=a+bs,\qquad a,b\in\mathbf Z_p.
 \tag{S33.5}
\]

\par\medskip\noindent\textbf{2. The self-dual trace-symplectic normal form}\quad 

For completeness, the precise normal form used in the coefficient calculation is justified here. The Rosati condition says \(\theta(ax,y)=\theta(x,ay)\) for \(a\in K\). Nondegeneracy of the field trace determines a unique alternating \(K\)-bilinear form \(\beta\) on \(L\otimes\mathbf Q_p\) by \[
 \operatorname{Tr}_{K/\mathbf Q_p}(a\beta(x,y))=\theta(ax,y)
 \quad(a\in K).
\] The balance identity proves bilinearity in both variables. Alternation follows from \(\theta(ax,x)=0\), using the balance identity and skew symmetry in characteristic zero.

For \(x,y\in L\), integrality after multiplying \(x\) by every \(a\in O\) gives \(\beta(x,y)\in O^\vee=s^{-1}O\), where \(O^\vee\) is the trace-dual lattice. The trace pairing identifies \(\operatorname{Hom}_O(L,s^{-1}O)\) with \(\operatorname{Hom}_{\mathbf Z_p}(L,\mathbf Z_p)\). Since the local form \(\theta\) is perfect, \(\beta\) is a perfect \(s^{-1}O\)-valued alternating form. Thus \(\psi=s\beta\) is a unimodular alternating \(O\)-bilinear form.

Choose a symplectic \(O\)-basis \((e_1,f_1,e_2,f_2)\) for \(\psi\). One obtains it by choosing a primitive vector, a second vector pairing to a unit, rescaling that vector to make the pairing one, and splitting off the orthogonal hyperbolic plane; the residual rank-two summand is handled in the same way. No division by two is involved. In this basis \[
 L=U_1\oplus U_2,
 \quad U_i=Oe_i\oplus Of_i,
 \quad \psi(e_i,f_j)=\delta_{ij},
 \quad \theta=\operatorname{Tr}_{K/\mathbf Q_p}(\psi/s).
 \tag{S33.6}
\] In particular each \(U_i\) is a unimodular symplectic \(\mathbf Z_p\)-lattice of rank four for the restricted trace form.

\par\medskip\noindent\textbf{3. The mixed second-Chern-class coefficients}\quad 

The degree-four part of (S33.3) gives \[
 \operatorname{ch}_2(\mathcal E)=\frac{\ell^2}{2N}-uC.
\] Since \(\ell^2/2\) is integral on a torus and \(c_2(\mathcal E)\) is integral, the identity \(\operatorname{ch}_2=\ell^2/2-c_2\) shows \[
 uC-\frac{\ell^2}{2N}\in H^4(X,\mathbf Z_p).
 \tag{S33.7}
\] For \(x,y\in O\), evaluate this four-form on \((xe_1,f_1,ye_2,f_2)\). The two embedding components of the trace form give \[
 C(xe_1,f_1,ye_2,f_2)=\operatorname{Tr}_{K/\mathbf Q_p}(xy/d),
\] while the orthogonality of \(U_1,U_2\) gives \[
 (\ell^2/2)(xe_1,f_1,ye_2,f_2)
 =\operatorname{Tr}(rx/s)\operatorname{Tr}(ry/s).
 \tag{S33.8}
\] Indeed each embedding component on one \(O\)-hyperbolic plane has rank two, so its square vanishes there; the mixed term in \(\theta_\sigma^2/2\) is the product of its two plane restrictions, with coefficient one.

In the basis \((1,s)\) of \(O\), the trace Gram matrix in (S33.8) for \(C\) is \(\operatorname{diag}(2/d,2)\). The corresponding linear functional for \(\ell\) has values \(2b,2a\). Therefore (S33.7) implies \[
 \frac{b^2}{N}-\frac{u}{2d}\in\mathbf Z_p,
 \qquad \frac{a^2}{N}\in\mathbf Z_p,
 \qquad \frac{ab}{N}\in\mathbf Z_p.
 \tag{S33.9}
\] Here division by four is legitimate because \(p\) is odd. These displayed factors are tied to the convention \(\theta=\operatorname{Tr}(\psi/s)\); no \(1/(2s)\) normalization has been substituted.

Let \(m=v_p(N)\) and \(k=v_p(b)\). The first congruence forces \(b\ne0\) and \[
 m=2k+1>0,
 \qquad v_p(a)\ge k+1,
 \qquad v_\pi(r)=2k+1=m.
 \tag{S33.10}
\] The summands \(a\) and \(bs\) have distinct \(\pi\)-valuations, at least \(2k+2\) and exactly \(2k+1\), respectively, proving the last equality. Moreover \[
 \frac{r^2}{N}
 =\frac{a^2}{N}+d\frac{b^2}{N}
             +2\frac{ab}{N}s
 \equiv\frac u2\pmod\pi.
 \tag{S33.11}
\] The first summand has positive \(p\)-valuation, the second has the asserted residue by (S33.9), and the third is divisible by \(\pi\). In particular \(u-r^2/N\) is a unit of \(O\).

\par\medskip\noindent\textbf{4. The integral Euler characteristic gives the contradiction}\quad 

The tangent bundle of an abelian variety is trivial, so Riemann--Roch (Stacks Project, Section~42.66, tag \texttt{02UO}) gives \(\chi(\mathcal E)=\int_X\operatorname{ch}_4(\mathcal E)\in\mathbf Z\). The character (S33.3) yields \[
 \operatorname{ch}_4(\mathcal E)
 =\frac{\ell^4}{24N^3}
       -\frac{u\ell^2C}{2N^2}+\frac{u^2}{N}\nu.
\] Using (S33.1) and writing \(r_\pm=\sigma_\pm(r)\), one computes \[
 \int_X\frac{\ell^4}{24}=D(r_+r_-)^2,
 \qquad
 \int_X\ell^2C=2D(r_+^2+r_-^2).
\] It follows that \[
 \chi(\mathcal E)
 =D\frac{\operatorname{Nm}(r)^2-Nu\operatorname{Tr}(r^2)+N^2u^2}{N^3}
 =D\frac{\operatorname{Nm}_{F/\mathbf Q}(Nu-r^2)}{N^3}.
 \tag{S33.12}
\] By (S33.11), \(Nu-r^2=N(u-r^2/N)\) has \(\pi\)-valuation \(2m\). Its local field norm has \(p\)-valuation \(2m\), since this ramified quadratic extension has residue degree one. Therefore \[
 v_p(\chi(\mathcal E))=2m-3m=-m<0,
 \tag{S33.13}
\] contradicting the integer Euler characteristic. This proves the main assertion.

For the semiregularity corollary, S-22 gives \(4D(Nq)^2/|G|=k\in\{2,4\}\). If an odd ramified prime does not divide \(D|G|\), then \(u=Nq\) is a \(p\)-adic unit. The degree-two character equation forces \(B=-\ell/N\), so the proposed normalized character is exactly (S33.3). The Néron--Severi hypothesis gives \(\ell\in F\theta\). The contradiction already proved shows that every odd ramified prime divides \(D|G|\). At \(2\), when \(2\nmid D|G|\), the value \(k=2\) would give the rational square \(u^2\) odd \(2\)-adic valuation \(-1\), which is impossible; \(k=4\) gives a \(2\)-adic unit, to which S-32 applies. This proves the combined condition (S33.4).
\end{proof}
\begin{remark}[An independent third-character certificate]
There is also a local contradiction before using the top-degree Euler calculation. This records the exact coefficient and its sign.

Let \(\theta_{\sigma,i}\) be the embedding component of \(\theta|_{U_i}\). Put \[
 \eta_1=\theta_{+,1}\theta_{-,1}
       =\theta|_{U_1}^2/2.
\] This is a primitive \(\mathbf Z_p\)-volume form on \(U_1\). In the particular ordered basis \((e_1,se_1,f_1,sf_1)\) its coefficient is \(4\) up to sign, which is harmless at odd \(p\). The degree \((4,2)\) component for the direct sum \(U_1\oplus U_2\) satisfies \[
 \begin{aligned}
 (\ell^3/6)_{(4,2)}
    &=\operatorname{Nm}(r)\eta_1\ell|_{U_2},\\
 (\ell C)_{(4,2)}
    &=\eta_1\theta_{\bar r}|_{U_2}.
 \end{aligned}
\] The second formula follows by writing \(C=\theta_{+,1}\theta_{+,2}+\theta_{-,1}\theta_{-,2}\): wedging its first summand with \(\ell|_{U_1}\) picks out \(r_-\), and wedging its second picks out \(r_+\). Thus \[
 \bigl(\operatorname{ch}_3(\mathcal E)\bigr)_{(4,2)}
 =\eta_1\theta_S|_{U_2},
 \qquad
 S=\frac{\operatorname{Nm}(r)r}{N^2}-\frac{u\bar r}{N}
   =\frac{\bar r(r^2-Nu)}{N^2}.
 \tag{S33.14}
\] For any integral degree-two form \(\ell\) on a torus, all divided powers \(\ell^j/j!\) are integral: in an exterior basis the coefficient of each ordered \(2j\)-monomial is a Pfaffian with integral entries. Hence \[
 \operatorname{ch}_3(\mathcal E)
 =\frac{\ell^3}{6}-\frac{\ell c_2(\mathcal E)}2
                         +\frac{c_3(\mathcal E)}2
 \in H^6(X,\mathbf Z_p)
\] for every odd \(p\), including \(p=3\). Since \(\eta_1\) is primitive and the restricted local trace form is perfect, (S33.14) would imply \(S\in O\). But (S33.10)--(S33.11) imply \[
 v_\pi(S)=v_\pi(\bar r)-v_\pi(N)
                  +v_\pi(r^2/N-u)=m-2m+0=-m<0.
\] This independently contradicts integrality. The argument uses only the usual integral Chern classes of a perfect object and the exterior algebra of a torus; no general integral topological Chern-character theorem is needed.
\end{remark}

\begin{theorem}[S-34: Determinant descent excludes every rank-one equivariant balanced seed]
\label{thm:semireg-rm-rank-one-equivariant}
Let $(X,\theta)$ be a polarized complex abelian fourfold
with Rosati-invariant real multiplication by a genuine
real quadratic field $F$, of type $(2,2)$. Write
$\theta=\theta_1+\theta_2$ and
\[
 v_q=\left(1-\frac q2\theta_1^2\right)
       \left(1-\frac q2\theta_2^2\right),\qquad q\in\mathbf Q_{>0}.
 \tag{S34.1}
\]
Let $G\subset X\times\widehat X$ be a finite effective
subgroup with a specified coherent action on
$\operatorname{Perf}(X)$ by translation and tensor product
with degree-zero line bundles. Let $E$ be a perfect complex
of virtual rank one satisfying
$e^B\operatorname{ch}(E)=v_q$ for a rational algebraic
divisor class $B$. No lift $\widetilde E$ to
$\operatorname{Perf}(X)^G$ which is simple and gluable
in that category is fully semiregular.

Here simplicity means
$\operatorname{End}_{\operatorname{Perf}(X)^G}
(\widetilde E)=\mathbf C$, and gluability means the
vanishing of its negative self-Ext groups. Neither condition
is required of the underlying object $E$. No maximal-order
or minimal-divisor-space hypothesis is imposed.
\end{theorem}
\begin{proof}
Write the action functors and their coherent multiplication
as
\[
 \phi_g(K)=t_{a_g}^*K\otimes P_g,
 \qquad
 \mu_{g,h}:\phi_g\phi_h\xrightarrow{\sim}\phi_{gh}.
\]
A lift consists of maps
$\lambda_g:\phi_g(E)\xrightarrow{\sim}E$ satisfying
\[
 \lambda_{gh}\circ\mu_{g,h,E}
       =\lambda_g\circ\phi_g(\lambda_h).
 \tag{S34.2}
\]
All natural identifications below use these specified
multiplication maps; replacing them by arbitrary
isomorphisms of the same functors would lose the argument.

Put $L=\det(E)$. The determinant of a perfect complex and
its compatibility with pullback are supplied by
Knudsen--Mumford, \emph{The projectivity of the moduli
space of stable curves I}, Theorem~2, p.~42, with the
extended determinant construction on pp.~40--42.
For a line bundle $P$ and a perfect complex of virtual
rank $r$, the local alternating determinant construction
gives
\[
 \det(E\otimes P)=\det(E)\otimes P^r,
 \qquad \det(c\operatorname{id}_E)=c^r\operatorname{id}_{\det E}.
 \tag{S34.3}
\]
Since $r=1$, taking determinants in (S34.2) gives maps
\[
 \delta_g:t_{a_g}^*L\otimes P_g\xrightarrow{\sim}L
 \tag{S34.4}
\]
with exactly the original scalar coherence, not a nontrivial
power of it. More explicitly, composition of (S34.4) for
$h$ and then $g$ agrees with (S34.4) for $gh$ after the
line-bundle identification induced by $\mu_{g,h}$.
This follows by applying the determinant functor to
(S34.2); the scalar part of $\mu_{g,h}$ is unchanged by
(S34.3). The formula can also be checked on any local
bounded complex of vector bundles, where the exponent
is the alternating sum of its ranks, namely one.

Conjugate the action by $K\mapsto K\otimes L^{-1}$.
Its functor for $g$ becomes
\[
 K\longmapsto t_{a_g}^*K\otimes R_g,
 \qquad
 R_g=t_{a_g}^*L\otimes P_g\otimes L^{-1}.
 \tag{S34.5}
\]
The maps $\delta_g$ trivialize $R_g$, and the coherence
just proved says that these trivializations identify the
entire action with the ordinary coherent translation action
of $G$ through $g\mapsto a_g$. They cancel any scalar
two-cocycle as well as the tensor factors. In particular
$E'=E\otimes L^{-1}$ acquires an ordinary translation
linearization. If $a_g=0$, (S34.4) implies $P_g\cong
\mathcal O_X$. Effectivity then gives $g=1$. Thus
$G\to X$ is injective. Its image $H$ acts freely, and
\[
 \pi:X\longrightarrow Y=X/H
 \tag{S34.6}
\]
is a finite etale isogeny of degree $|G|$.

For clarity, perfect descent here can be realized directly,
without assuming that descent in an arbitrary derived
homotopy category is effective. On $\pi_*E'$ the ordinary
linearization gives a genuine finite-group action. Form its
averaging idempotent and its image
\[
 E_Y=\operatorname{im}\left(
  \frac1{|G|}\sum_{g\in G}g:\pi_*E'\longrightarrow\pi_*E'
                         \right).
 \tag{S34.7}
\]
The pushforward is perfect: etale locally on $Y$ the
cover is a disjoint union of copies of $Y$, and the
assertion is then the preservation of perfection by a
finite direct sum. Perfection is fpqc local by Stacks
Project, Lemma~36.12.4, tag \texttt{09UG}.
The perfect category is closed under direct summands, so
the idempotent has a perfect image. After pulling back to
$X$, flat base change identifies $\pi^*\pi_*E'$ with
$\bigoplus_{h\in H}t_h^*E'$. The linearization identifies
the invariant summand with the diagonal copy of $E'$.
Consequently $\pi^*E_Y\cong E'$, compatibly with descent.
The same construction on morphism complexes, or the
adjunction and exactness of finite-group invariants, gives
for every integer $j$
\[
 \operatorname{Ext}^j_Y(E_Y,E_Y)
      \cong \operatorname{Ext}^j_X(E',E')^G
      \cong
      \operatorname{Ext}^j_{\operatorname{Perf}(X)^G}
                   (\widetilde E,\widetilde E).
 \tag{S34.8}
\]
Thus the descended object is simple and gluable. The
underlying sheaf descent used by this construction is
Stacks Project, Proposition~35.5.2, tag \texttt{023T};
the explicit finite averaging construction supplies the
perfect-complex statement and the Ext comparison.

Full semiregularity is also preserved in this descent.
One may use the enhanced equivalence
$\operatorname{Perf}(Y)\simeq
\operatorname{Perf}(X)^H$ just constructed and naturality
of the full Hochschild trace, as in S-10 and Perry 2026,
Remark~2.3. Equivalently, suppose
$u\in\operatorname{Ext}^2_Y(E_Y,E_Y)$ has zero full trace.
Its pullback has zero full trace on $X$. Under (S34.8)
it is an invariant class. The forgetting map on
$HH_{-2}$ is injective for this translation action by
the identity-sector argument of S-21. Naturality
therefore makes the equivariant full trace of $u$ zero;
full semiregularity of the lift forces $u=0$.
Tensoring by $L^{-1}$ is an equivalence and preserves
this property. Hence a fully semiregular lift would
give a fully semiregular rank-one perfect $E_Y$ on $Y$.

The character and the real-multiplication data descend
with the required normalization. The degree-two component
of $e^B\operatorname{ch}(E)=v_q$ gives $B=-c_1(E)$,
so $\operatorname{ch}(E')=v_q$.
Conjugating the rational $F$-action by the isogeny
defines an $F$-action on $Y$, and
$\theta_Y=(\pi^*)^{-1}\theta$ is a rational ample class
for which that action is Rosati invariant. Choose a
positive integer $m$ such that $\theta'=m\theta_Y$
is an integral polarization. With the corresponding
two real components of $\theta'$, the equality of
characters reads
\[
 \operatorname{ch}(E_Y)
   =\left(1-\frac{q/m^2}{2}(\theta'_1)^2\right)
      \left(1-\frac{q/m^2}{2}(\theta'_2)^2\right).
 \tag{S34.9}
\]
All equalities follow by pulling back along $\pi$,
which is an isomorphism on rational cohomology. The
polarization volume changes by
$\int_Y(\theta')^4/24=m^4D/|G|$, and $q/m^2>0$.
S-30 excludes a simple gluable fully semiregular
rank-one object with (S34.9), independently of the
polarization degree or an integral maximal-order action.
This is the desired contradiction.
\end{proof}
